%---------------------------------------------------------------------------------------------------------------------------- 
\graphicspath{{Parti/P03-Meccanica_SR/C02-Statica_CR/Figure/}}
\chapter{Statica del corpo rigido}\label{cap:statica_CR}
%---------------------------------------------------------------------------------------------------------------------------- 


\begin{sommario}
	\todo{Da rivedere}Si richiamano le equazioni di equilibrio del corpo rigido
	come motivazione per la riduzione di un sistema di forze
	alle sue grandezze globali. Si definiscono risultante e
	momento risultante, si introduce la coppia di forze e si
	deriva la legge di variazione del momento al variare del
	polo. Si studiano gli invarianti del campo dei momenti e
	si introduce l'asse centrale del sistema di forze con una
	costruzione geometrica parallela a quella del
	\CAPref{cap:cinematica_infinitesima}. Si classificano i
	sistemi equivalenti e si specializza la trattazione al
	caso piano, con la determinazione del centro di riduzione.
	Si conclude con la riduzione di sistemi di forze
	equiorientate e il calcolo del baricentro delle forze.
\end{sommario}

\section{La nozione di forza e di lavoro}\label{sec:forza_lavoro}

Un corpo $\Cc$  è generalmente immerso in un ambiente più vasto, con cui scambia
azioni meccaniche. Occuparsi della statica di $\Cc$ significa rendere conto di \emph{questa
	interazione fra il corpo e ciò che sta al di fuori di esso}. 

 Per farlo, ci sono essenzialmente due strade.
La prima  interpreta l'interazione in termini vettoriali: la si può pensare come un'azione di contatto o a distanza, 
assimilabile ad `spinta' o una `trazione' che il corpo riceve, in un punto o in una 
regione limitata del dominio che occupa, dotata di un'intensità, di
una direzione e di un verso, capace di
alterare lo stato della configurazione di $\Cc$. È la lettura che conduce naturalmente 
alla nozione di \emph{forza applicata}.

La seconda  interpreta l'interazione in termini energetici: la stessa azione, quando il
corpo si sposta, compie un lavoro, ovvero trasferisce  energia fra il
corpo e l'esterno lungo il moto realmente compiuto. Non è una
costruzione ausiliaria rispetto alla prima, ma la stessa interazione,
letta questa volta come una quantità scalare anziché come un vettore: 
il \emph{lavoro reale}.

Che le due letture --- vettoriale ed energetica --- siano davvero
equivalenti, e non due nozioni indipendenti, non è scontato a priori. Entrambe conducono a due formulazioni equivalenti della nozione di
equilibrio.

Per quanto la nostra intuizione sembri aver ormai in qualche modo introiettato la nozione di forza,
collegandola ad esempio al peso di un oggetto o alla spinta che possiamo imprimergli,
darne una definizione rigorosa e fisicamente coerente non è affatto banale.
Storicamente, la questione si è rivelata più insidiosa di quanto
sembri. Le prime formulazioni moderne restano legate a
un'idea sfuggente: la forza descritta da Leonardo, ad
esempio, è poco più di un'energia invisibile insita nei corpi, capace di
manifestarsi solo nei suoi effetti: ``una virtù spirituale, una potenza invisibile''.  Anche volendo appoggiarsi alla
dinamica, definendo la forza attraverso l'effetto che produce sul moto, come suggerito dalla meccanica Newtoniana,
ci si scontra con una circolarità di fondo: si misura la forza tramite
l'accelerazione che essa provoca, ma per farlo occorre già sapere cosa
sia una forza, per poterla eventualmente calibrare o misurare. Un
passo decisivo verso una nozione operativa arriva solo con
l'osservazione, dovuta a Hooke, che una forza può essere \emph{misurata}
attraverso l'allungamento che produce in un corpo elastico di
riferimento: non si definisce così cosa sia una forza, ma si costruisce
un modo riproducibile di confrontarne e sommarne gli effetti --- il che,
ai fini della statica, è ciò che davvero serve. Tuttavia è chiaro che la nozione di forza
rimane di fatto oscura, esistendo soltanto per il tramite dei suoi effetti. Il mondo visibile 
consta essenzialmente di attributi cinematici: è il moto che il nostro occhio percepisce,
è la deformazione che, in ultima analisi, il dinamometro misura, collegandola solo in seconda battuta 
alla forza stessa.

Ma la misura non è l'unico modo, né il più immediato, in cui una forza
si rende conoscibile. Se si vuole sapere quanto pesa un oggetto, prima
ancora di posarlo su una bilancia, lo si \textit{soppesa}: gli si
imprime uno spostamento e si valuta la `fatica' che questo costa. Il nostro
corpo non sa misurare direttamente una forza, ma sa valutare il lavoro
che essa compie. Questa seconda via, apparentemente meno immediata da formalizzare ma
non meno fondamentale, apre al cosiddetto
\textit{metodo dei lavori virtuali}.



\medskip 

Seguiamo qui la stessa via già scelta per il corpo (\SECref{sez.CR}): non
definiamo la forza, ma la caratterizziamo attraverso le sue proprietà
formali e i suoi effetti --- una volta nella forma di un vettore
applicato, un'altra volta nella forma di un funzionale lineare sugli
spostamenti (virtuali) che il ragionamento precedente ci ha condotto a
introdurre. È a questa duplice caratterizzazione che dedichiamo il resto
della sezione.



\subsection{Due modi di rappresentare l'azione meccanica}

\paragraph{La forza come vettore applicato.}


Il primo modo è quello già introdotto in generale per un singolo vettore
applicati: una forza agente in un punto
$P$  è una coppia $(P,\fb)\in \Ec\times\Uc$. 

Ricordando la \eqref{eq:momento_polare}, chiamiamo momento della forza $\fb$ applicata in $P$, rispetto al polo $O$, la quantità
\begin{equation}
		\bm{\upmu}(O) =\overrightarrow{OP} \times \fb.
\end{equation}


\paragraph{Il lavoro di una forza.}
Per formalizzare la seconda strada abbiamo bisogno di un concetto nuovo, quello di lavoro.
Il lavoro compiuto da una forza $\fb$ applicata in $P$, quando questo abbia uno spostamento $\ub$ è la
quantità scalare
\begin{equation}
\Wc(\fb)[\ub]:=\fb\cdot\ub.
\end{equation}
Per $\fb$ fissato, $\Wc$ definisce dunque un'applicazione $\Wc: \Uc\to\Ro$, che associa a ogni spostamento $\ub$ il lavoro compiuto dalla forza; la notazione
con le parentesi quadre indica che l'applicazione $\Wc$, \textit{valutata} su un certo campo vettoriale $\ub$, agisce linearmente su $\ub$.

Questa definizione prescinde dalla circostanza che $\ub$ sia quello che pertiene un corpo rigido. Specializziamola al caso in cui $\ub=\ub_{\rm rig}\in\Uc_{\rm rig}$, dove per $\ub_{\rm rig}$
intendiamo un campo di spostamento rispettoso della condizione \eqref{eq:campo_rigido}:
\begin{equation}
\begin{aligned}
		\Wc(\fb)[\ub_{\rm rig}]&=\fb(P)\cdot (\ub(O)+\bm{\upvartheta}\times\overrightarrow{OP})\\
		&=\fb(P)\cdot \ub(O)+\overrightarrow{OP}\times \fb\cdot \bm{\upvartheta}\\
		&=\fb(P)\cdot \ub(O)+\bm\upmu(O)\cdot\bm{\upvartheta},
\end{aligned}
	\end{equation}
avendo sfruttato la proprietà del prodotto misto
$\ab \cdot (\bb \times \cb) = (\cb \times \ab) \cdot \bb$ e la definizione di momento $\bm\upmu(O)$ della forza $\fb$ rispetto al polo $O$.

Osserviamo che  alle due quantità \textit{cinematiche} rilevanti in uno spostamento rigido $\ub(O)$ e $\bm\upvartheta$, 
la nozione di lavoro associa altrettanti descrittori  \textit{dinamici}\footnote{Qui le accezioni `cinematica'  e `dinamica' sono da intendersi in senso etimologico: la prima attiene al moto (\textgreek{k\'{i}nhma} -- \textit{kìnema}), la seconda alla forza (\textgreek{d\'{u}namis} -- \textit{d\'{y}namis}) . }: la forza e il suo momento. 
Il lavoro dunque istituisce una forma di \textit{dualità}: la forza è duale alla componente traslazionale del moto, il momento della forza alla componente rotazionale, nel senso che $\fb$ compie lavoro su $\ub$, mentre $\bm{\upmu}$ compie lavoro su $\bm{\upvartheta}$. 



Diremo che il lavoro è \textit{virtuale} quando il campo $\ub_{\rm rig}$ su
cui lo si valuta non è quello effettivamente realizzato da un moto del
corpo, ma un campo scelto arbitrariamente fra tutti quelli rigidi
ammissibili --- cioè uno qualunque dei campi $\ub(O)+\bm{\upvartheta}\times
\overrightarrow{OP}$, al variare di $\ub(O)$ e $\bm{\upvartheta}$, anche
quando il corpo non lo compie affatto. In altri termini, la posizione
descritta da $\ub_{\rm rig}$ non è necessariamente quella che $\Cc$ occupa
realmente, ma una qualunque fra le infinite posizioni che $\Cc$, in linea
di principio, potrebbe occupare. Questo permette di concepire una serie
di \textit{Gedankenexperiment}, esperimenti mentali --- spostamenti ipotizzati, non necessariamente
realizzati --- rispetto ai quali valutare il lavoro delle forze applicate
a $\Cc$.

\subsection{Sistemi di forze e coppie}



\paragraph{Sistemi di forze.}
Consideriamo un sistema $\mathcal{S}$ di $n$ forze applicate:
$\Sc=\{(P_k,\,\fb_k), \;k=1,\dots,n\}$. 

Definiamo \textit{risultante}\index{risultante|textbf}\nomenclature[L]{$\fb$}{risultante di un sistema di forze (vettore libero); forza} è il vettore libero
\begin{equation}\label{eq:risultante}
	\rb \coloneqq \sum_{k=1}^{n} \fb_k.
\end{equation}
%di componenti $X = \sum_k X_k$, $Y = \sum_k Y_k$,
%$Z = \sum_k Z_k$.

Il \textit{momento risultante}\index{momento!risultante|textbf}\nomenclature[G]{$\bm{\upmu}(O)$}{momento risultante rispetto al polo $O$ (vettore assiale)} rispetto a un polo $O$ è il
vettore assiale
\begin{equation}\label{eq:momento_risultante}
	\bm{\upmu}(O) \coloneqq \sum_{k=1}^{n}
	\overrightarrow{OP_k} \times \fb_k,
\end{equation}
somma dei momenti\index{momento!di una forza} $\bm{\upmu}_k(O) =
\overrightarrow{OP_k} \times \fb_k$ dei singoli vettori
applicati.
% (\FIGref{fig:momento_forza}\,a). 

%In componenti
%(\FIGref{fig:momento_forza}\,b):
%\begin{equation}\label{eq:momento_risultante_det}
%	\bm{\upmu}(O) = \sum_{k=1}^{n}
%	\begin{vmatrix}
%		\ab_x & \ab_y & \ab_z \\
%		x_k   & y_k   & z_k   \\
%		X_k   & Y_k   & Z_k
%	\end{vmatrix}
%	= \sum_{k=1}^{n}
%	\begin{bmatrix}
%		0    & -z_k &  y_k \\
%		z_k  &  0   & -x_k \\
%		-y_k &  x_k &  0
%	\end{bmatrix}
%	\begin{Bmatrix} X_k \\ Y_k \\ Z_k \end{Bmatrix},
%\end{equation}
%dove $x_k$, $y_k$, $z_k$ sono le coordinate di $P_k$
%rispetto a $O$.

%
%\begin{figure}[!htbp]
%	\centering
%	\def\ax{0.5}\def\ay{0.30}
%	\providecommand{\xyz}[3]{{#2 + (#1)*\ax},{#3 + (#1)*\ay}}
%	\providecommand{\rightangle}[3]{%
%		\coordinate (ra1) at ($(#2)!2.4mm!(#1)$);
%		\coordinate (ra3) at ($(#2)!2.4mm!(#3)$);
%		\draw[line width=0.5pt] (ra1) -- ($(ra1)+(ra3)-(#2)$) -- (ra3);}
%	\begin{minipage}[b]{0.46\textwidth}\centering
%		\begin{tikzpicture}[>=Stealth, line width=0.8pt, scale=1.0]
%			% ===== PANNELLO (a) =====
%			\draw[gray!60] (\xyz{-1.7}{-1.2}{0}) -- (\xyz{-1.7}{3.50}{0})-- (\xyz{3.0}{3.50}{0}) -- (\xyz{3.0}{-1.2}{0}) -- cycle;
%			\node[gray!75] at (\xyz{-1.2}{-1.0}{0}) {$\pi$};
%			\coordinate (O)  at (\xyz{0}{0}{0});
%			\coordinate (Pk) at (\xyz{0.5}{2.0}{0});
%			\coordinate (Ff) at (\xyz{1.7}{2.3}{0});
%			\coordinate (Of) at (\xyz{1.2}{0.3}{0});
%			\coordinate (MuB) at (\xyz{0}{0}{1.25});
%			\coordinate (MuC) at (\xyz{0}{0}{1.47});
%			\draw[gray!55, thin] (O) -- (Of) -- (Ff) -- (Pk) -- cycle;
%			
%			%				\draw[->, thin] (O) -- (Pk);
%			%				\node[font=\small, below] at ($(O)!0.5!(Pk)$) {$\overrightarrow{OP_k}$};
%			
%			\draw[->,thick] (O) -- (Pk) node[font=\small, midway, sloped, below] {$\overrightarrow{OP_k}$};
%			
%			\draw[->, thick] (Pk) -- (Ff) node[right=-2pt] {$\fb_k$};
%			\draw[->>, thick] (O) -- (MuB) node[left=1pt, yshift=-10pt] {$\bm{\upmu}_k(O)$};
%			%\draw[->, thin] ([shift={(0.18,0)}]MuC)arc[x radius=0.18, y radius=0.07, start angle=0, end angle=305];
%			\begin{scope}[gray!55]
%				\rightangle{MuB}{O}{Of}   % angolo retto fra mu e la retta // f_k
%				\rightangle{MuB}{O}{Pk}
%			\end{scope}
%			\fill (O)  circle (1.5pt) node[left=2pt] {$O$};
%			\fill (Pk) circle (1.5pt) node[below right=-1pt] {$P_k$};
%		\end{tikzpicture}\\[0.8ex]
%		(a)
%	\end{minipage}\hfill
%	\begin{minipage}[b]{0.50\textwidth}\centering
%		\begin{tikzpicture}[>=Stealth, line width=0.8pt, scale=1.0]
%			% proiezione: x in avanti (basso-sinistra), y destra, z su
%			\def\ax{-0.65}\def\ay{-0.48}
%			
%			\def\az{0.7}                            % <1 = comprime in altezza
%			\renewcommand{\xyz}[3]{{#2 + (#1)*\ax},{(#3)*\az + (#1)*\ay}}
%			% componenti di P_k
%			\def\xk{1.4}\def\yk{2.6}\def\zk{2.0}
%			\coordinate (O)   at (\xyz{0}{0}{0});
%			\coordinate (Px)  at (\xyz{\xk}{0}{0});
%			\coordinate (Py)  at (\xyz{0}{\yk}{0});
%			\coordinate (Pxy) at (\xyz{\xk}{\yk}{0});
%			\coordinate (Pk)  at (\xyz{\xk}{\yk}{\zk});
%			
%			% --- scatola di decomposizione (base + verticale) ---
%			\draw[gray!55, thin] (O) -- (Px) -- (Pxy) -- (Py) -- cycle;
%			\draw[gray!55, thin] (Pxy) -- (Pk);
%			
%			% --- assi x,y,z ---
%			\draw[->, thin] (O) -- (\xyz{2.6}{0}{0}) node[below left=-2pt] {$x$};
%			\draw[->, thin] (O) -- (\xyz{0}{3.4}{0}) node[right]          {$y$};
%			\draw[->, thin] (O) -- (\xyz{0}{0}{2.9}) node[above]          {$z$};
%			
%			% --- vettore posizione OP_k ---
%			%				\draw[->, thin] (O) -- (Pk);
%			%				\node[font=\small, above left=-8pt, yshift=6pt] at ($(O)!0.42!(Pk)$) {$\overrightarrow{OP_k}$};
%			
%			% --- terna delle componenti di forza in P_k ---
%			\coordinate (Xt) at (\xyz{\xk+0.85}{\yk}{\zk});
%			\coordinate (Yt) at (\xyz{\xk}{\yk+0.85}{\zk});
%			\coordinate (Zt) at (\xyz{\xk}{\yk}{\zk+0.85});
%			\draw[->, thick] (Pk) -- (Xt) node[above=4pt] {$X_k$};
%			\draw[->, thick] (Pk) -- (Yt) node[right=-1pt] {$Y_k$};
%			\draw[->, thick] (Pk) -- (Zt) node[above=-1pt] {$Z_k$};
%			
%			% --- etichette componenti (spigoli O->Py->Pxy->Pk) ---
%			\node[font=\small, below=20pt, xshift=-20pt]      at ($(O)!0.5!(Py)$)   {$y_k$};
%			\node[font=\small,  right=1pt, yshift=-6pt] at ($(Py)!0.5!(Pxy)$) {$x_k$};
%			\node[font=\small, right, yshift=6pt]      at ($(Pxy)!0.5!(Pk)$) {$z_k$};
%			
%			% --- punti ---
%			\fill (O)  circle (1.5pt) node[above left=-1pt]  {$O$};
%			\fill (Pk) circle (1.5pt) node[above right=-1pt] {$P_k$};
%			
%			% ===== MOMENTI in O (interni, assiali ->> + ricciolo) =====
%			\coordinate (Mx) at (\xyz{1.2}{0}{0});
%			\coordinate (My) at (\xyz{0}{1.2}{0});
%			\coordinate (Mz) at (\xyz{0}{0}{1.4});
%			\draw[->>, thick] (O) -- (Mx);
%			%\draw[->, thin, rotate around={126:(Mx)}]([shift={(0.16,0)}]Mx) arc[x radius=0.16, y radius=0.06, start angle=0, end angle=305];
%			\node[font=\small, above,  xshift=-4pt, yshift=6pt] at (Mx) {$\bm{\upmu}_{xk}(O)$};
%			\draw[->>, thick] (O) -- (My);
%			%\draw[->, thin, rotate around={-90:(My)}] ([shift={(0.16,0)}]My) arc[x radius=0.16, y radius=0.06, start angle=0, end angle=305];
%			\node[font=\small, below=-1pt, xshift=-6pt] at (My) {$\bm{\upmu}_{yk}(O)$};
%			\draw[->>, thick] (O) -- (Mz);
%			%\draw[->, thin, rotate around={0:(Mz)}]	([shift={(0.16,0)}]Mz) arc[x radius=0.16, y radius=0.06, start angle=0, end angle=305];
%			\node[font=\small, left=1pt] at (Mz) {$\bm{\upmu}_{zk}(O)$};
%			
%			% ===== FORZE lungo gli assi (esterne, -> ; traslabili con \sx,\sy,\sz) =====
%			\def\Lm{0.9}
%			\def\sx{1.2}
%			\coordinate (Xa) at (\xyz{\sx}{0}{0});
%			\coordinate (Xb) at (\xyz{\sx+\Lm}{0}{0});
%			\draw[->, thick] (Xa) -- (Xb);
%			\node[font=\small, below=3pt] at (\xyz{\sx+\Lm/2}{0}{0}) {$X_k$};
%			\def\sy{2.1}
%			\coordinate (Ya) at (\xyz{0}{\sy}{0});
%			\coordinate (Yb) at (\xyz{0}{\sy+\Lm}{0});
%			\draw[->, thick] (Ya) -- (Yb);
%			\node[font=\small, above] at (Yb) {$Y_k$};
%			\def\sz{1.5}
%			\coordinate (Za) at (\xyz{0}{0}{\sz});
%			\coordinate (Zb) at (\xyz{0}{0}{\sz+\Lm});
%			\draw[->, thick] (Za) -- (Zb);
%			\node[font=\small, left=2pt] at (Zb) {$Z_k$};
%			
%		\end{tikzpicture}
%		\\[0.8ex]
%		(b)
%	\end{minipage}
%	\caption{(a)~Momento $\bm{\upmu}_k(O)$ rispetto al polo
%		$O$ della forza $(P_k,\fb_k)$, definito come
%		$\bm{\upmu}^O_k = \protect\overrightarrow{OP_k}\times\fb_k$.
%		(b)~Componenti cartesiane $x_k,y_k,z_k$ del vettore
%		posizione $\protect\overrightarrow{OP_k}$ e componenti
%		$X_k,Y_k,Z_k$ della forza $\fb_k$.}
%	\label{fig:momento_forza}
%\end{figure}

Il risultante e il momento risultante $(\fb,\,\bm{\upmu}(O))$ costituiscono gli
\textit{elementi caratteristici} del sistema $\Sc$ rispetto al
polo $O$.

\paragraph{Nozione di coppia.}
Oltre alla forza, un ente dinamico che gioca un ruolo fondamentale nelle applicazioni è la \textit{coppia}. 
Per motivarne l'introduzione, consideriamo un particolare sistema, composto da due forze eguali ed opposte, applicate
in due punti distinti:
\begin{equation}
	\Sc= \{ (P_1, \fb_1), (P_2, -\fb_1)  \}.
\end{equation}


Il risultante è ovviamente nullo. Il momento risultante rispetto a un polo $O$ qualsiasi è
\begin{equation}\label{eq:coppia}
	\bm{\upmu}(O)
	= \overrightarrow{OP_1} \times \fb
	+ \overrightarrow{OP_2} \times (-\fb)
	= (\overrightarrow{OP_1} - \overrightarrow{OP_2})
	\times \fb_1
	= \overrightarrow{P_2P_1} \times \fb_1,
\end{equation}
e \textit{non dipende dal polo} (\FIGref{fig:coppia}). Il suo modulo è
$|\bm{\upmu}| = |\fb_1|\,b$, dove
$b = |\protect\overrightarrow{P_2P_1}|\sin\alpha$ è il
\textit{braccio}\index{braccio}, ovvero la distanza fra le linee d'azione;
la direzione è perpendicolare al piano di azione della coppia  di forze e il verso
è quello della rotazione che la coppia tende a produrre,
secondo la regola della mano destra.

Sull'insieme delle coppie di forze --- coppie di vettori applicati
$(P_1,\fb_1)$, $(P_2,-\fb_1)$ --- introduciamo la relazione che dichiara
equivalenti due coppie con lo stesso momento risultante $\bm{\upmu}$.
Si tratta di una relazione di equivalenza, indotta semplicemente
dall'uguaglianza fra vettori $\bm{\upmu}$: per quanto visto, $\bm{\upmu}$
è infatti un invariante completo, nel senso che due coppie di forze sono
equivalenti secondo questa relazione se e solo se producono lo stesso
$\bm{\upmu}$. Chiamiamo \textit{coppia}\index{coppia|textbf} $\cb$ ciascuna
classe di equivalenza così definita: ogni coppia di forze $(P_1,\fb_1)$,
$(P_2,-\fb_1)$ con un dato momento risultante $\bm{\upmu}$ è allora una
\textit{realizzazione} di $\cb$, e tutte le realizzazioni di $\cb$ sono,
ai fini dell'equilibrio del corpo $\Cc$, indistinguibili.



\begin{figure}[htbp]
	\centering
	\newcommand{\xyz}[3]{{#2 + (#1)*\ax},{#3 + (#1)*\ay}}
	\newcommand{\rightangle}[3]{%
		\coordinate (ra1) at ($(#2)!2.4mm!(#1)$);
		\coordinate (ra3) at ($(#2)!2.4mm!(#3)$);
		\draw[line width=0.5pt] (ra1) -- ($(ra1)+(ra3)-(#2)$) -- (ra3);}
	\tikzset{midarrow/.style={postaction={decorate},
			decoration={markings, mark=at position #1 with {\arrow{Stealth}}}},
		midarrow/.default=0.55}
	\begin{tikzpicture}[>=Stealth, line width=0.8pt, scale=1.2]
		\def\ax{0.5}\def\ay{0.30}
		% piano pi
		%\draw[gray!60] (\xyz{-1.2}{-2.4}{0}) -- (\xyz{-1.2}{4.8}{0}) -- (\xyz{3.8}{4.8}{0}) -- (\xyz{3.8}{-2.4}{0}) -- cycle;
		\draw[gray!60] (\xyz{-1.60}{-2.4}{0}) -- (\xyz{-1.60}{3.8}{0})	-- (\xyz{4.8}{3.8}{0}) -- (\xyz{4.8}{-2.4}{0}) -- cycle;
		%\node[gray!75] at (\xyz{-1.85}{-2.2}{0}) {$\pi$};   % sposta anche l'etichetta pi
		\node[gray!75] at (\xyz{-1.2}{-2.0}{0}) {$\pi$};
		
		\def\Lf{1.3}
		\coordinate (O)  at (\xyz{2.8}{0.9}{0});
		\coordinate (P1) at (\xyz{2.0}{0.3}{0});
		\coordinate (P2) at (\xyz{0.6}{2.0}{0});
		\coordinate (F1) at (\xyz{2.0}{0.3-\Lf}{0});
		\coordinate (F2) at (\xyz{0.6}{2.0+\Lf}{0});
		
		% rette d'azione (grigie, prolungate)
		\draw[gray!55, thin] (\xyz{2.0}{-2.2}{0}) -- (\xyz{2.0}{1.0}{0});
		\draw[gray!55, thin] (\xyz{0.6}{-2.2}{0}) -- (\xyz{0.6}{.5}{0});
		% prolungamento grigio di P2P1 oltre P1
		\coordinate (PPe) at ($(P1)!-0.6!(P2)$);
		\draw[gray!55, thin] (P1) -- (PPe);
		
		% vettori posizione
		%		\draw[thin, midarrow] (O) -- (P1);
		%		\node[font=\small, above=1pt] at ($(O)!0.5!(P1)$) {$\overrightarrow{OP_1}$};
		
		\draw[thin, midarrow] (O) -- (P1) node[font=\small, midway, sloped, above] {$\overrightarrow{OP_1}$};
		
		\draw[thin, midarrow] (O) -- (P2);
		\node[font=\small, above right=-2pt] at ($(O)!0.5!(P2)-(0,0.2)$) {$\overrightarrow{OP_2}$};
		
		
		% P2P1
		%		\draw[thin, midarrow] (P2) -- (P1);
		%		\node[font=\small, below=1pt, xshift=-16pt, yshift=6pt] at ($(P2)!0.42!(P1)$) {$\overrightarrow{P_2P_1}$};
		
		\draw[thin, midarrow] (P2) -- (P1) node[font=\small, midway, sloped, below] {$\overrightarrow{P_2P_1}$};
		
		
		% forze
		\draw[->, line width=1.2pt] (P1) -- (F1) node[below left=-4pt] {$\fb_1$};
		\draw[->, line width=1.2pt] (P2) -- (F2) node[below, xshift=-16pt] {$-\fb_1$};
		
		% angolo alpha fra il prolungamento grigio di P2P1 e f_1, in P1
		\pic[draw, thick, angle radius=6mm, angle eccentricity=1.5,	"\small$\alpha$"] {angle = PPe--P1--F1};
		
		% braccio b sui prolungamenti grigi (perpendicolare fra le due rette)
		\coordinate (B1) at (\xyz{2.0}{-1.9}{0});   % su retta d'azione di f_1
		\coordinate (B2) at (\xyz{0.6}{-1.9}{0});   % su retta d'azione di -f_1
		\draw[<->, thin] (B1) -- (B2);
		\node[font=\small, above=1pt, xshift=-2pt, yshift=-5pt] at ($(B1)!0.5!(B2)$) {$b$};
		
		\begin{scope}[gray!55]
			\rightangle{B2}{B1}{\xyz{2.0}{-1.4}{0}}
			\rightangle{B1}{B2}{\xyz{0.6}{-1.4}{0}}
		\end{scope}
		
		
		
		% ===== momento mu LIBERO, a destra del piano, con DUE angoli retti =====
		\def\my{0.60}                      % <- alza/abbassa cambiando solo questo
		\coordinate (Mb) at (3.6,\my);
		\coordinate (Mt) at ($(Mb)+(0,1.4)$);
		\coordinate (Mc) at ($(Mb)+(0,1.62)$);
		\coordinate (pax) at ($(Mb)+(0.55,0.33)$);  % direzione tipo a_x
		\coordinate (pay) at ($(Mb)+(0.95,0)$);     % direzione tipo a_y
		\draw[gray!55, thin] (Mb) -- (pax) -- ($(pax)+(pay)-(Mb)$) -- (pay) -- cycle;
		\draw[->>, thick] (Mb) -- (Mt) node[left=1pt] {$\bm{\upmu}$};
		\draw[->, thin] ([shift={(0.18,0)}]Mc)	arc[x radius=0.18, y radius=0.07, start angle=0, end angle=305];
		\begin{scope}[gray!55]
			\rightangle{Mt}{Mb}{pax}
			\rightangle{Mt}{Mb}{pay}
		\end{scope}
		
		% punti
		\fill (O)  circle (1.4pt) node[above=2pt] {$O$};
		\fill (P1) circle (1.4pt) node[above left=-2pt] {$P_1$};
		\fill (P2) circle (1.4pt) node[below] {$P_2$};
	\end{tikzpicture}
	\caption{Momento della coppia generata dalle forze $\fb_1$ e
		$\fb_2=-\fb_1$ applicate in $P_1$ e $P_2$. Il momento
		$\bm{\upmu}=\protect\overrightarrow{P_2P_1}\times\fb_1$ è perpendicolare
		al piano delle forze, con modulo $|\fb_1|\,b$ e braccio
		$b=|\protect\overrightarrow{P_2P_1}|\sin\alpha$, ed è indipendente dal
		polo $O$.}
	\label{fig:coppia}
\end{figure}


\paragraph{Sistemi di forze e coppie.}
Un sistema di forze e coppie è un insieme di $n_f$ forze e $n_c$ coppie:
\begin{equation}
\Sc=\{(P_k,\,\fb_k) \in \Ec\times\Uc, \;k=1,\dots,n_f; \;\; \cb_i\in\Uc, \; j=1, \dots, n_c\}
\end{equation}

In tal caso, il risultante e il momento risultante rispetto ad un polo $O$ si definiscono come
\begin{equation}\label{eq:riduzione_generalizzata}
	\rb = \sum_{k=1}^{n_f} \fb_k, \qquad
	\bm{\upmu}(O) = \sum_{k=1}^{n_f}
	\overrightarrow{OP_k} \times \fb_k
	+ \sum_{j=1}^{n_c} \cb_j.
\end{equation}
Le coppie contribuiscono al momento risultante ma non al
risultante.

\paragraph{Densità di forze.}
In molte applicazioni le forze agenti su un corpo non sono
concentrate in punti isolati ma distribuite su volumi,
superfici o linee. Si introducono le corrispondenti
\textit{densità di forze}\index{densita di forze@densità di forze|textbf}: la densità di volume $\fb_V$
(dimensioni fisiche $[F\,L^{-3}]$), la densità di superficie
$\fb_S$ ($[F\,L^{-2}]$) e la densità di linea $\fb_L$
($[F\,L^{-1}]$).

Il risultante e il momento risultante si ottengono
sostituendo le sommatorie con integrali sui rispettivi domini:
\begin{equation}\label{eq:risultante_distribuite}
	\rb = \int_V \fb_V \dd V
	+ \int_S \fb_S \dd S
	+ \int_L \fb_L \dd L,
\end{equation}
\begin{equation}\label{eq:momento_distribuite}
	\bm{\upmu}(O) = \int_V \overrightarrow{OP}
	\times \fb_V \dd V
	+ \int_S \overrightarrow{OP}
	\times \fb_S \dd S
	+ \int_L \overrightarrow{OP}
	\times \fb_L \dd L,
\end{equation}
dove $\overrightarrow{OP}$ è il vettore posizione del punto
generico $P$ rispetto al polo $O$. La legge di
variazione~\eqref{eq:variazione_momento}, gli invarianti e
tutte le proprietà del campo dei momenti restano validi: la
coppia $(\fb,\,\bm{\upmu}(O))$ caratterizza il sistema
indipendentemente dal fatto che le forze siano concentrate o
distribuite.

\begin{oss}
	Il caso più frequente nella meccanica delle strutture è
	quello di forze distribuite lungo una linea (l'asse di una
	trave) o su una superficie (la faccia di una piastra o di
	un guscio). In questi casi la densità di forze prende il
	nome di \textit{carico distribuito}\index{carico!distribuito}: il carico di linea
	$\fb_L$ ha le dimensioni di una forza per unità di
	lunghezza, il carico di superficie $\fb_S$ di una forza
	per unità di area. Esempi tipici sono il peso proprio di
	una trave (carico di linea uniforme) e la pressione del
	vento su una parete (carico di superficie).
\end{oss}


\subsection{Il lavoro di un sistema di forze e coppie}
Definiamo innanzitutto il lavoro di una coppia, ripartendo dalla coppia di forze definite dal sistema
$\Sc=\{(P_1, \fb_1), (P_2, -\fb_1)\}$. Il lavoro computo dal sistema sarà la somma dei lavori di ogni forza:
\begin{equation}
\begin{aligned}
		\Wc(\Sc)[\ub_{\rm rig}]=\fb_1\cdot\ub(P_1)-\fb_1\cdot \ub(P_2)&=\fb_1\cdot \bm{\upvartheta}\times(\overrightarrow{OP_1}-\overrightarrow{OP_2})\\
		&=\overrightarrow{P_2P_1}\times \fb_1\cdot\bm{\upvartheta}=\bm{\upmu}\cdot \bm{\upvartheta}.
\end{aligned}
\end{equation}
Per traslato, il lavoro compiuto su uno spostamento rigido da una coppia $\cb$ è definito da
\begin{equation}
	\Wc(\cb)[\ub_{\rm rig}]\coloneqq \cb\cdot \bm{\upvartheta}.
\end{equation}

Possiamo quindi determinare il lavoro compiuto su uno spostamento rigido da un sistema di forze e coppie $\Sc=\{(P_k,\,\fb_k) \in \Ec\times\Uc, \;k=1,\dots,n_f; \;\; \cb_i\in\Uc, \; j=1, \dots, n_c\}$:
\begin{equation}
\begin{aligned}\label{LavSis}
	\Wc(\Sc)[\ub_{\rm rig}]&=\sum_{k=1}^{n_f}\fb_k\cdot(\ub(O)+\bm\upvartheta\times \overrightarrow{OP_i})+\sum_{j=1}^{n_c}\cb_i\cdot \bm{\upvartheta}\\
	&=\left(\sum_{k=1}^{n_f} \fb_k\right)\cdot\ub(O)+\left( \sum_{k=1}^{n_f} \overrightarrow{OP_i}\times\fb_i  +\sum_{j=1}^{n_c}\cb_i\right)\cdot \bm{\upvartheta}\\
	&=\rb\cdot\ub(O)+\bm{\upmu}(O)\cdot \bm{\upvartheta},
\end{aligned}
\end{equation}
dove $\big(\rb,\bm{\upmu}(O)\big)$ sono il risultante e il momento risultante del sistema $\Sc$ considerato.

Analogamente a quanto osservato per il caso della singola forza, il lavoro istituisce una dualità tra le variabili cinematiche 
che parametrizzano la classe degli spostamenti rigidi  e le variabili dinamiche \textit{globali} di un sistema di forze e coppie:
il risultante $\fb$ e il momento risultante $\bm{\upmu}(O)$ sono rispettivamente coniugati alla componente traslatoria $\ub(O)$ e rotatoria $\bm{\upvartheta}$.

\paragraph{Il lavoro delle densità di forza.}
Per un carico distribuito $\fb_L(P)$ lungo una linea, il
lavoro  è definito come
\begin{equation}
	\Wc = \int_L \fb_L(P) \cdot \ub(P)\,\dd L.
\end{equation}
Sostituendo la formula fondamentale
$\ub(P) = \ub(O) + \bm{\upvartheta} \times
\overrightarrow{OP}$ e procedendo come nel caso concentrato:
\begin{equation}
	\Wc = \rb \cdot \ub(O)
	+ \bm{\upmu}(O) \cdot \bm{\upvartheta},
\end{equation}
con
\begin{equation}
	\rb = \int_L \fb_L\,\dd L, \qquad
	\bm{\upmu}(O) = \int_L \overrightarrow{OP}
	\times \fb_L\,\dd L.
\end{equation}
La forma è identica al caso concentrato: la distinzione
fra forze concentrate e distribuite scompare a livello
di grandezze globali.

\section{L'equilibrio}

Avendo fornito una doppia caratterizzazione dell'azione meccanica che 
il mondo esterno esercita sul corpo $\Cc$, è naturale attendersi che 
le condizioni che garantiscono l'equilibrio di $\Cc$ possano essere date in modo duplice.
Entrambe si basano su due \textit{princip\^{\i}}: inerenti gli elementi caratteristici del sistema di forze $\Sc$
cui $\Cc$ è pensato essere assoggettato nel primo caso, il lavoro compiuto da $\Sc$ su un qualunque 
ammissibile spostamento di $\Cc$ stesso. 

\paragraph{Il principio di equilibrio.}
\textit{In una configurazione di equilibrio, il sistema di forze e coppie applicate sul corpo $\Cc$ deve essere \textit{nullo}, ovvero}
\begin{equation}\label{ECS}
\rb=\mathbf{0}, \qquad \bm{\upmu}(O)=\mathbf{0}.
\end{equation}

\medskip

Osserviamo che, benché il momento risultante dipenda dal polo, 
dalla~\eqref{eq:momento_risultante}
si trova:
\begin{equation}\label{eq:variazione_momento}
	\begin{aligned}
		\bm{\upmu}(P)
	&= \sum_k \overrightarrow{PP_k} \times \fb_k
	= \sum_k \overrightarrow{OP_k} \times \fb_k
	- \overrightarrow{OP} \times \sum_k \fb_k\\
	&= \bm{\upmu}(O) + \rb \times \overrightarrow{OP}.
	\end{aligned}
\end{equation}
Dunque, se il risultante $\rb$ è nullo come richiede il principio di equilibrio, il momento è invariante rispetto alla scelta del polo.

Le \eqref{ECS} sono anche note come \textit{Equazioni Cardinali della Statica}.

Come caso particolare notiamo infine che un sistema di due forze è
equilibrato se e solo se le due forze hanno la stessa retta
d'azione, uguale intensità e verso opposto --- costituiscono
cioè una coppia di braccio nullo; un sistema di tre forze non
parallele è equilibrato se e solo se le tre forze sono
complanari, concorrenti in un punto e hanno risultante nullo.
La coplanarità è conseguenza necessaria: la terza forza,
dovendo essere uguale e contraria al risultante delle prime
due, giace nel piano da esse individuato. La concorrenza
garantisce che il momento risultante rispetto al punto di
concorso sia nullo, essendo nullo il braccio di ciascuna
forza; poiché anche il risultante è nullo, il momento è nullo
rispetto a qualsiasi polo.


\paragraph{Il principio dei lavori virtuali.}
\textit{In una configurazione di equilibrio, il sistema di forze e coppie $\Sc$ applicate sul corpo $\Cc$ compie lavoro virtuale nullo
per ogni possibile moto rigido del corpo:}
	\begin{equation}\label{eq:PLV}
	\Wc(\Sc)[\ub_{\rm rig}] = \rb \cdot \ub(O)
	+ \bm{\upmu}(O) \cdot \bm{\upvartheta} = 0
	\qquad \forall\,\ub(O),\,\bm{\upvartheta}\in \Uc.
\end{equation}



\paragraph{Equivalenza delle due formulazioni.}
La dimostrazione dell'equivalenza delle due formulazioni è piuttosto immediata. 
Che il Principio di Equilibrio implichi il Principio dei Lavori Virtuali (PLV) è banale:
se $\rb=\bm{\upmu}=\mathbf{0}$, stante la \eqref{LavSis}, il lavoro è evidentemente nullo, qualunque sia $\ub(O)$ e $\bm{\upvartheta}$.

Per dimostrare il viceversa, utilizziamo la quantificazione che richiede l'annullarsi del 
lavoro virtuale \textit{per ogni} possibile moto, ovvero per ogni $\ub(O)$ e
$\bm{\upvartheta}$. Siamo dunque autorizzati a scegliere $\ub(O)\neq \mathbf{0}$ e
$\bm{\upvartheta}=\mathbf{0}$, ottenendo che $\Wc=\rb\cdot \ub(O)=0$ per ogni $\ub(O)$. 
Data l'arbitrarietà di  $\ub(O)$, questa condizione impone che $\rb$ debba essere ortogonale
a tutti i vettori di $\Uc$ e l'unica possibilità che questo accada è per il vettore nullo, sicché $\rb=\mathbf{0}$.
In maniera analoga, prendendo $\ub(O)= \mathbf{0}$ e
$\bm{\upvartheta}\neq\mathbf{0}$, si conclude che $ \bm{\upmu}(O)=\mathbf{0}$.


\begin{oss}
Come specificato, le due formulazioni rappresentano  dei \textit{postulati} e, in quanto tali, si accettano come veri senza bisogno di dimostrazione. 
In molti trattati si assume vero solo il postulato di equilibrio, mentre si deriva --- allo stesso modo che noi abbiamo impiegato
per mostrare l'equivalenza dei due princip\^{\i} --- la condizione sul lavoro virtuale. In tal caso, si parla più propriamente di \textit{Teorema dei Lavori Virtuali}. Molti trattatisti,
con una certa libertà, sovrappongo Principio e Teorema.
\end{oss}

\begin{oss}
Alle due equazioni cardinali della statica ci si riferisce generalmente come equilibrio alla \textit{traslazione} e alla \textit{rotazione}, rispettivamente.
La ragione di questa dizione comunemente utilizzata può essere compresa solo se si introduce la nozione di lavoro, che  istituisce la dualità tra 
le due variabili dinamiche presenti nelle equazioni e le variabili cinematiche che ne giustificano la qualificazione.

\end{oss}



\section{Invarianza rispetto al cambiamento di osservatore}\label{sec:invarianza_osservatore}

L'equilibrio di $\Cc$ è stato finora caratterizzato in due modi equivalenti: l'annullamento di risultante e
momento risultante, oppure l'annullamento del lavoro
virtuale su ogni moto rigido (PLV). C'è un terzo modo di arrivare alla
stessa conclusione, che  parte da un requisito di
\textit{oggettività}: l'azione meccanica di un sistema di forze esterne su
$\Cc$ è un fatto fisico, e il lavoro che la rappresenta non può dipendere
da quale, fra osservatori ugualmente legittimi, ne valuta l'entità.

\paragraph{Lavoro di un sistema e cambiamento di osservatore.}

Per un sistema di forze e coppie con elementi caratteristici
$(\rb,\bm{\upmu}(O))$, il lavoro è dato dalla \eqref{LavSis}. 
Per costruzione, $\Wc$ è un funzionale \textit{lineare} di
$(\ub(O), \bm{\upvartheta})$: non ha termine
costante, ed è additivo, $\Wc(\Sc)[\ub_1+\ub_2]=\Wc(\Sc)[\ub_1]+\Wc(\Sc)[\ub_2]$.

Un cambiamento di osservatore roto-traslazionale è descritto da un moto
rigido ausiliario arbitrario $\ub^+\in\Uc_{\rm rig}$, parametrizzato dalla coppia $(\ub^+(O),\bm{\upvartheta}^+)\in\Uc\times\Uc$: è la deriva relativa fra i due osservatori. Se il primo
osservatore attribuisce a $\Cc$ un moto virtuale $\ub$, il secondo,
essendo esso stesso in moto rigido $\ub^+$ rispetto al primo,
attribuisce al medesimo corpo il moto apparente $\ub+\ub^+$, e
calcola di conseguenza
\begin{equation}\label{eq:lavoro_osservatore_cambiato}
	\Wc^+(\Sc) := \Wc(\Sc)[\ub+\ub^+], \qquad \ub, \ub^+\in\Uc_{\rm rig}.
\end{equation}
Le quantità $\rb$ e $\bm{\upmu}(O)$ non cambiano: sono il dato fisico del
problema, lo stesso per entrambi gli osservatori; a cambiare è solo il
moto che ciascuno attribuisce al corpo.

\paragraph{Il postulato di invarianza.}

Un sistema di forze esterne $\Sc$ agenti su $\Cc$ soddisfa il postulato di
	invarianza se il lavoro virtuale che gli si attribuisce non dipende dal
	cambiamento di osservatore:
	\begin{equation}\label{eq:postulato_invarianza}
		\Wc(\Sc)[\ub] = \Wc(\Sc)[\ub+\ub^+]
		\qquad \forall\,\ub,\ub^+\in\Uc_{\rm rig}.
	\end{equation}


Il contenuto fisico della \eqref{eq:postulato_invarianza} è che il lavoro
misurato non può registrare la deriva $\ub^+$ di un osservatore
rispetto a un altro: se lo facesse, $\Wc$ misurerebbe qualcosa
sull'osservatore, non sul sistema di forze.

Vale  il seguente teorema di invarianza rispetto all'osservatore:

\medskip

\index{invarianza rispetto all'osservatore!equivalenza con PLV|textbf}
	\emph{Un sistema di forze e coppie soddisfa il postulato di invarianza
	\eqref{eq:postulato_invarianza} se e solo se soddisfa il principio dei
	lavori virtuali, ovvero se e solo se $\rb=\zerovb$ e $\bm{\upmu}(O)=\zerovb$.}
	
	\medskip



	Per la linearità di $\Wc(\Sc)$,
	\[
	\Wc(\Sc)[\ub+\ub^+] = \Wc(\Sc)[\ub] + \Wc(\Sc)[\ub^+],
	\]
	quindi la \eqref{eq:postulato_invarianza} equivale a
	\begin{equation}\label{eq:invarianza_plv}
		\Wc(\Sc)[\ub^+] = 0 \qquad \forall\,\ub^+\in\Uc_{\rm rig},
	\end{equation}
	indipendentemente dallo spostamento $\ub$ di partenza: è esattamente il PLV.
	Valutando la \eqref{eq:invarianza_plv} su $\ub^+$ puramente
	traslatorio ($\bm{\upvartheta}^+=\zerovb$, $\ub^+(O)$
	arbitrario) si ottiene $\rb\cdot\ub^+(O)=0$ per ogni
	$\ub^+(O)\in\Uc$, dunque $\rb=\zerovb$; su $\ub^+$
	puramente rotazionale si ottiene analogamente
	$\bm{\upmu}(O)=\zerovb$. Il viceversa è immediato: se
	$\rb=\bm{\upmu}(O)=\zerovb$, allora $\Wc(\Sc)\equiv 0$, e la
	\eqref{eq:postulato_invarianza} vale banalmente per ogni $\ub$ e
	ogni $\ub^+\in\Uc_{\rm rig}$.

\section{Struttura dei sistemi di forze}\label{sec:struttura_sistemi_forze}

Le formulazioni dell'equilibrio viste finora ---  il principio di equilibrio, il
PLV, l'invarianza rispetto all'osservatore --- rispondono tutte alla stessa
domanda: \textit{quando} un sistema di forze e coppie sia equilibrato.
Presuppongono però di conoscere già $(\rb,\bm{\upmu}(O))$ rispetto a un
polo comodo. Restano aperte alcune domande complementari, che riguardano
non l'equilibrio ma la \textit{struttura} di un sistema di forze,
equilibrato o no: quali invarianti possiedano gli elementi caratteristici, quando due sistemi diversi
producano lo stesso effetto meccanico, quali forme essenziali un sistema
di forze possa assumere. Sono gli strumenti a cui dedichiamo il resto del
capitolo, e torneranno utili più avanti.



\subsection{Invarianti del campo dei momenti}\label{sec:invarianti_F}
Ricordiamo che il momento varia rispetto al polo come prescritto dalla \eqref{eq:trasporto_momento}, ovvero:

\begin{equation}\label{eq:variazione_momento}
\bm{\upmu}(P)
		= \bm{\upmu}(O) + \rb \times \overrightarrow{OP}.
\end{equation}


Prima di procedere, conviene mettere a confronto la struttura di questo
campo con quella già incontrata nel Cap~\ref{cap:cinematica_infinitesima}:
i risultati ottenuti lì per il campo di spostamento --- dipendendo solo
dalla forma algebrica del campo, non dal suo significato fisico ---
valgono qui verbatim, previa la sostituzione indicata in tabella.
Questo evita di ripetere calcoli identici nelle sezioni seguenti.

\begin{center}
	\renewcommand{\arraystretch}{1.4}
	\begin{tabular}{ccc}
		\toprule
		\textbf{Campo di spostamento} & & \textbf{Campo dei momenti} \\
		\midrule
		spostamento $\ub(P)$ & & momento $\bm{\upmu}(P)$ \\
		spostamento al polo: $\ub(O)$ & & momento al polo: $\bm{\upmu}(O)$ \\
		rotazione $\bm{\upvartheta}$ & & risultante $\rb$ \\
		invariante $\bm{\upvartheta}\cdot\ub(O)$ & & invariante $\rb\cdot\bm{\upmu}(O)$ \\
		asse centrale degli spostamenti & & asse centrale delle forze \\
		distanza $|\ub_\perp(O)|/|\bm{\upvartheta}|$ & & distanza $|\bm{\upmu}_\perp(O)|/|\rb|$ \\
		\bottomrule
	\end{tabular}
\end{center}

Il campo dei momenti possiede dunque due invarianti, indipendenti dal
polo:
\begin{enumerate}[label=(\roman*)]
	\item la risultante $\rb$ (dato del sistema, non del polo);
	\item il prodotto scalare $\rb\cdot\bm\upmu(O)$ --- l'equiproiettività
	del campo dei momenti, per la tabella sopra e lo stesso conto di
	sottosez.~\ref{sez:interpretazione}: se $P$ è allineato con $O$ nella
	direzione di $\rb$, l'intero vettore $\bm\upmu$ --- non solo la sua
	componente scalare --- è invariante.
\end{enumerate}


\subsection{Asse centrale del sistema di forze}

Per la tabella sopra, e per la stessa costruzione geometrica dell'asse centrale degli spostamenti rigidi (\SECref{assecentraleSP}), esiste, per $\rb\neq\zerovb$, un'unica
retta parallela a $\rb$ --- l'\textit{asse centrale del sistema di
	forze}\index{asse centrale!del sistema di forze|textbf} --- lungo la
quale il momento è puramente parallelo a $\rb$:
\begin{equation}\label{eq:asse_centrale_forze}
	\overrightarrow{OC}
	= \frac{\rb\times\bm\upmu(O)}{|\rb|^2}
	+ \lambda\,\frac{\rb}{|\rb|}, \qquad \lambda\in\Ro.
\end{equation}
Lungo l'asse centrale il momento vale $\bm\upmu_\parallel
= (\rb\cdot\bm\upmu(O))/|\rb|^2\;\rb$; la distanza dell'asse centrale dal
polo $O$ è $|\bm\upmu_\perp(O)|/|\rb|$; nel caso di forza unica
($\rb\cdot\bm\upmu(O)=0$) l'asse centrale è il luogo dei poli a momento
nullo, cioè la retta d'azione della risultante
(\FIGref{f:asse_centrale_forze}).

\begin{figure}[htbp]
	\centering
	\usetikzlibrary{arrows.meta,calc,decorations.markings,angles,quotes}
	% --- proiezione obliqua: (x=profondita', y=oriz, z=vert) -> schermo ---
	\def\ax{0.5}\def\ay{0.30}
	\newcommand{\xyz}[3]{{#2 + (#1)*\ax},{#3 + (#1)*\ay}}
	\tikzset{midarrow/.style={postaction={decorate},decoration={markings, mark=at position 0.56 with {\arrow{Stealth}}}}}
	\begin{tikzpicture}[scale=1.4, >=Stealth, line width=0.8pt]
		
		%==== coordinate 3D (x=profondita', y=oriz, z=vert) ====
		\coordinate (Pi1) at (\xyz{-1.3}{-1.0}{0});   % era {-1.3}{-0.7}
		\coordinate (Pi2) at (\xyz{-1.3}{ 4.5}{0});   % era {-1.3}{ 3.5}
		\coordinate (Pi3) at (\xyz{ 2.7}{ 4.5}{0});   % era { 2.7}{ 3.5}
		\coordinate (Pi4) at (\xyz{ 2.7}{-1.0}{0});   % era { 2.7}{ 3.5}→ refuso: era { 2.7}{ 3.5}
		
		% O e suo asse
		\coordinate (O)     at (\xyz{1.6}{0.5}{0});
		\coordinate (Otop)  at (\xyz{1.6}{0.5}{2.7});
		\coordinate (Obot)  at (\xyz{1.6}{0.5}{-0.8});
		\coordinate (ThOa)  at (\xyz{1.6}{0.5}{1.10});
		\coordinate (ThOb)  at (\xyz{1.6}{0.5}{1.90});
		\coordinate (curlO) at (\xyz{1.6}{0.5}{2.12});
		\coordinate (uparO) at (\xyz{1.6}{0.5}{0.90});
		
		% u_perp(O): ora verso il BASSO-SINISTRA
		\coordinate (uperpO) at (\xyz{0.76}{0.26}{0});  % era {0.2}{0.1}
		
		% P_pi davanti-destra (OP_pi va in basso-destra): base 0.2,2.3 (era 1.886,2.5)
		\coordinate (Ppi)   at (\xyz{0.2}{2.3}{0});
		\coordinate (Ctop)  at (\xyz{0.2}{2.3}{2.7});
		\coordinate (Cbot)  at (\xyz{0.2}{2.3}{-0.8});
		\coordinate (ThCa)  at (\xyz{0.2}{2.3}{1.10});
		\coordinate (ThCb)  at (\xyz{0.2}{2.3}{1.90});
		\coordinate (curlC) at (\xyz{0.2}{2.3}{2.12});
		\coordinate (uparC) at (\xyz{0.2}{2.3}{0.90});
		
		% prodotto vettoriale = -u_perp(O): ora verso l'ALTO-DESTRA
		\coordinate (cross)  at (\xyz{1.04}{2.54}{0});  % era {1.6}{2.7}
		\coordinate (Bpar) at (\xyz{-0.64}{2.06}{0});
		
		
		%==== piano pi ====
		\draw[gray!60] (Pi1) -- (Pi2) -- (Pi3) -- (Pi4) -- cycle;
		\node[gray!75] at ($(Pi1)+(0.58,0.14)$) {$\pi$};
		
		%==== assi verticali ====
		\draw[gray!55, thin] (Obot) -- (Otop);
		\draw[gray!55, thin] (Cbot) -- (Ctop);
		
		%==== parallelogramma di lati u_perp(O) e OP_pi (evidenziato) ====
		\filldraw[gray!10, draw=gray!55, thin]	(O) -- (uperpO) -- (Bpar) -- (Ppi) -- cycle;
		
		%==== triangolo O--C_perp--(f x OC_perp) ====
		\fill[gray!10] (O) -- (Ppi) -- (cross) -- cycle;
		
		% u_perp(O) traslato in P_pi: opposto e collineare al prodotto vettoriale
		\draw[->>,  thin] (Ppi) -- (Bpar);
		%\node[font=\small, below=4pt, right] at (Bpar) {$\ub_{\perp}(O)$};
		
		%==== segmento O--(punta di vartheta x OP_pi) e ampiezza di rotazione vartheta ====
		\draw[draw=gray!75, thin] (O) -- (cross);
		%\pic[draw=black, ->, thin, angle radius=10mm, angle eccentricity=1.25,	"$\vartheta$"{font=\small, inner sep=1pt}] {angle = Ppi--O--cross};
		
		%==== vettori rotazione assiali + riccioli (sopra la doppia freccia) ====
		\draw[->, thick] (ThOa) -- (ThOb) node[right=1pt] {$\rb$};
		\draw[->, thick] (ThCa) -- (ThCb) node[right=1pt] {$\rb$};
		
		%==== costruzione nel piano ====
		%\draw[->, thin] (O) -- (Ppi);
		%\draw[midarrow, thin] (O) -- (Ppi);   % era  \draw[->, thin] (O) -- (Ppi);
		%\node[font=\small, below=1pt] at ($(O)!0.30!(Ppi)$) {$\overrightarrow{OC_\perp}$};   % era 0.55, below=2pt
		
		\draw[midarrow,thin] (O) -- (Ppi) node[midway, sloped, below,xshift=-12pt] {$\overrightarrow{OC_\perp}$};
		
		
		\draw[->>, thick] (O) -- (uperpO);
		\node[above=1pt] at (uperpO) {$\bm{\upmu}_{\perp}(O)$};
		\draw[->>, thick] (O) -- (uparO);
		\node[right=1pt] at ($(O)!0.6!(uparO)$) {$\bm{\upmu}_{\parallel}$};
		
		%==== a P_pi ====
		\draw[->>, thick] (Ppi) -- (uparC);
		\node[right=1pt] at ($(Ppi)!0.6!(uparC)$) {$\bm{\upmu}_{\parallel}$};
		\draw[->>, thick] (Ppi) -- (cross);
		\node[below right=-6pt] at (cross) {$\rb\times\overrightarrow{OC_\perp}$};
		
		
		
		%==== punti ====
		\fill (O)   circle (1.5pt) node[above right=-1pt] {$O$};
		\fill (Ppi) circle (1.5pt) node[below right=-4pt] {$C_\perp$};
		
		%==== annotazione asse centrale ====
		\node[align=left, font=\footnotesize] (lab) at ($(Ctop)+(1.8,-0.3)$)	{asse centrale della\\ sollecitazione};
		\draw[->, thin] ($(lab.north west)+(0,-0.18)$) -- ($(Ctop)!0.18!(ThCb)$);
		
	\end{tikzpicture}%
	\caption{Costruzione geometrica dell'asse centrale del
		sistema di forze. Il momento rispetto a $O$ si decompone
		nella componente $\bm{\upmu}^O_{\parallel}$, parallela a
		$\rb$ e invariante per tutti i poli, e nella componente
		$\bm{\upmu}^O_{\perp}$, giacente nel piano $\pi$
		ortogonale a $\rb$. Il punto $C_\perp$, intersezione
		dell'asse centrale con $\pi$, è determinato dalla
		condizione $\rb \times \protect\overrightarrow{OC_\perp}
		= -\bm{\upmu}^O_{\perp}$; l'asse centrale è la retta
		per $C_\perp$ parallela a $\rb$.}
	\label{f:asse_centrale_forze}
\end{figure}


L'asse centrale del sistema di forze gode delle seguenti
proprietà:
\begin{enumerate}[label=(\roman*)]
	\item per $\rb \neq \zerovb$ è univocamente determinato e
	la sua posizione è indipendente dal polo $O$ usato per la
	costruzione;
	\item lungo di esso il momento è puramente parallelo a
	$\rb$ e pari a
	$\bm{\upmu}_{\parallel}
	= (\rb \cdot \bm{\upmu}(O))/|\rb|^2\;\rb$;
	\item la distanza dell'asse centrale dal polo $O$ è
	$|\bm{\upmu}_{\perp}(O)|/|\rb|$;
	\item nel caso di forza unica
	($\rb \cdot \bm{\upmu}(O) = 0$), la componente invariante
	$\bm{\upmu}_{\parallel}$ è nulla e l'asse centrale è il
	luogo dei poli a momento nullo: è la retta d'azione della
	risultante.
\end{enumerate}

\begin{oss}
	La condizione che definisce l'asse centrale può essere
	espressa anche in forma algebrica: un punto $P$ appartiene
	all'asse centrale se e solo se
	$\bm{\upmu}(P) \times \rb = \zerovb$, cioè se il momento
	è parallelo alla risultante. Sostituendo
	la~\eqref{eq:variazione_momento} e sviluppando il doppio
	prodotto vettoriale si ottiene l'equazione
	\begin{equation}\label{eq:asse_centrale_algebrica_F}
		\bm{\upmu}(O) \times \rb
		+ |\rb|^2\,\overrightarrow{OP}
		- (\rb \cdot \overrightarrow{OP})\,\rb
		= \zerovb,
	\end{equation}
	che definisce la stessa retta costruita geometricamente
	nei paragrafi precedenti.
\end{oss}

In base ai valori degli invarianti si distinguono i seguenti
casi:
\begin{enumerate}[label=(\roman*)]
	\item \textit{Dinamo (o sistema avvitante)}\index{sistema di forze!dinamo (sistema avvitante)}:
	$\rb \neq \zerovb$,
	$\bm{\upmu}(O) \neq \zerovb$,
	$\rb \cdot \bm{\upmu}(O) \neq 0$. Il momento ha sia
	componente nel piano ortogonale a $\rb$ sia componente
	lungo $\rb$. Sull'asse centrale il momento è puramente
	parallelo a $\rb$: il sistema, visto da un polo sull'asse
	centrale, appare come una \textit{forza più una coppia
		parallela}.
	\item \textit{Forza unica}\index{sistema di forze!forza unica}:
	$\rb \neq \zerovb$ e
	$\rb \cdot \bm{\upmu}(O) = 0$. L'asse centrale è il luogo
	dei poli a momento nullo; il sistema è equivalente a una
	singola forza $\rb$ applicata su tale retta.
	\item \textit{Coppia pura}\index{coppia!pura}:
	$\rb = \zerovb$, $\bm{\upmu}(O) \neq \zerovb$. Il momento
	è indipendente dal polo; l'asse centrale non esiste (o,
	equivalentemente, è all'infinito).
	\item \textit{Sistema nullo}: $\rb = \zerovb$,
	$\bm{\upmu}(O) = \zerovb$. Nessun effetto statico.
\end{enumerate}

\begin{oss}
	La corrispondenza con la classificazione dei moti rigidi
	del \CAPref{cap:cinematica_infinitesima} è immediata: la dinamo corrisponde alla
	rototraslazione, la forza unica alla rotazione pura, la
	coppia pura alla traslazione pura, il sistema nullo al
	riposo. Si noti l'inversione apparente: nella cinematica è
	la rotazione pura ad annullare l'invariante scalare
	($\bm{\upvartheta} \cdot \ub(O) = 0$, spostamento nullo
	sull'asse centrale); nella statica è la forza unica
	($\rb \cdot \bm{\upmu}(O) = 0$, momento nullo sull'asse
	centrale). In entrambi i casi, l'asse centrale è il luogo
	dei punti in cui il campo si annulla.
\end{oss}

\todo[inline]{La parte sotto va un po' sistemata recuperando ciò che manca della vecchia versione, si tratta solo di spostare}
%------------------------------------------------------------
\subsection{Equivalenza di sistemi di forze}\label{sec:equivalenza}
%------------------------------------------------------------

Due sistemi di forze che producono lo stesso effetto su ogni
corpo rigido --- cioè la stessa risultante e lo stesso momento
risultante rispetto a un polo qualsiasi --- sono
\textit{equivalenti}\index{sistema di forze!equivalenza|textbf}\index{equivalenza|see{sistema di forze}}. Il concetto di equivalenza è il
fondamento delle operazioni di riduzione e semplificazione
dei sistemi di forze: sostituire un sistema con uno equivalente
più semplice non altera l'equilibrio del corpo.

\paragraph{Definizione di equivalenza.}

Due sistemi di forze $\mathcal{S}_1$ e $\mathcal{S}_2$ si
dicono equivalenti se valgono le condizioni
\begin{equation}\label{eq:def_equivalenza}
	\rb_{1} = \rb_{2}, \qquad
	\bm{\upmu}_{1}(O) = \bm{\upmu}_{2}(O).
\end{equation}
La legge di variazione del
momento~\eqref{eq:variazione_momento} garantisce che, se
le condizioni~\eqref{eq:def_equivalenza} valgono per un
polo $O$, allora valgono per ogni polo: è dunque
sufficiente verificarle rispetto a un solo polo.


\begin{oss}
	La relazione così definita è un'equivalenza in senso
	proprio: è riflessiva (ogni sistema è equivalente a sé
	stesso), simmetrica (se $\mathcal{S}_1$ è equivalente a
	$\mathcal{S}_2$, allora $\mathcal{S}_2$ è equivalente a
	$\mathcal{S}_1$) e transitiva (se $\mathcal{S}_1$ è
	equivalente a $\mathcal{S}_2$ e $\mathcal{S}_2$ è
	equivalente a $\mathcal{S}_3$, allora $\mathcal{S}_1$ è
	equivalente a $\mathcal{S}_3$). L'insieme di tutti i
	sistemi di forze risulta dunque partizionato in classi di
	equivalenza, ciascuna individuata da una coppia
	$(\rb, \bm{\upmu}(O))$.
\end{oss}

\begin{oss}
	L'equivalenza è una relazione fra \textit{sistemi}, non
	fra singole forze: due sistemi equivalenti producono lo
	stesso effetto statico su un corpo rigido, pur potendo
	essere composti da forze diverse per numero, intensità e
	punto di applicazione.
\end{oss}



\paragraph{Operazioni elementari di equivalenza.}

Le seguenti operazioni trasformano un sistema $\mathcal{S}$ in
un sistema $\mathcal{S}'$ equivalente. Esse costituiscono gli
strumenti fondamentali per la semplificazione e riduzione dei
sistemi di forze.

\bigskip 
%\paragraph{Traslazione lungo la linea d'azione.}
Una forza $\fb_k$ applicata in $P_k$ può essere spostata in un
qualsiasi altro punto $P_k'$ appartenente alla sua linea
d'azione senza alterare né la risultante né il momento
risultante (\FIGref{fig:operazioni_elementari_1}\,a). Infatti, poiché $\overrightarrow{P_k P_k'}$ è
parallelo a $\fb_k$:
\begin{equation}
	\overrightarrow{OP_k'} \times \fb_k
	= (\overrightarrow{OP_k}
	+ \overrightarrow{P_k P_k'}) \times \fb_k
	= \overrightarrow{OP_k} \times \fb_k.
\end{equation}

\bigskip

Due o più forze le cui rette d'azione si intersecano in un
punto $P$ possono essere traslate lungo le rispettive rette
d'azione fino a $P$ e lì sostituite dalla loro risultante (\FIGref{fig:operazioni_elementari_1}\,b). Se
$\fb_1$ e $\fb_2$ sono concorrenti in $P$:
\begin{equation}
	(P,\,\fb_1),\;(P,\,\fb_2)
	\quad\longleftrightarrow\quad
	(P,\,\fb_1+\fb_2).
\end{equation}
L'operazione inversa è la \textit{decomposizione}: una forza
può essere sostituita da due o più forze applicate nello stesso
punto la cui somma vettoriale sia uguale alla forza originale (\FIGref{fig:operazioni_elementari_1}\,c).

\begin{figure}[htbp]
	\centering
	\subfloat[][\emph{traslazione lungo la linea d'azione}]{%
		\begin{tikzpicture}[>=Stealth, line width=0.8pt, scale=0.80]
			% linea d'azione (direzione di fb_k)
			\draw[dashed] (0,0) -- (4.5,2.25);
			% punto Pk
			\fill (1,0.5) circle (1.5pt) node[below right] {$P_k$};
			% forza in Pk
			\draw[->, thick] (1,0.5) -- (2.6,1.3)
			node[above left] {$\fb_k$};
			% punto Pk' più avanti sulla stessa retta
			\fill (3,1.5) circle (1.5pt) node[below right] {$P_k'$};
			% forza in Pk'
			\draw[->, thick] (3,1.5) -- (4.6,2.3)
			node[above left] {$\fb_k$};
	\end{tikzpicture}}
	\qquad
	\subfloat[][\emph{composizione di forze concorrenti}]{%
		\begin{tikzpicture}[>=Stealth, line width=0.8pt, scale=0.80]
			% punto di concorrenza P
			\fill (2,1.5) circle (1.5pt) node[above right] {$P$};
			% linee d'azione
			\draw[dashed] (0,0) -- (2,1.5);
			\draw[dashed] (4,0) -- (2,1.5);
			% forze f1 e f2 dirette verso P
			\draw[->, thick] (0.6,0.45) -- (2,1.5)
			node[midway, left] {$\fb_1$};
			\draw[->, thick] (3.4,0.45) -- (2,1.5)
			node[midway, right] {$\fb_2$};
			% risultante in P
			\draw[->, thick] (2,1.5) -- (2,3.0)
			node[right] {$\fb_1+\fb_2$};
	\end{tikzpicture}}
	\qquad
	\subfloat[][\emph{decomposizione di una forza}]{%
		\begin{tikzpicture}[>=Stealth, line width=0.8pt, scale=0.80]
			% punto P
			\fill (1.5,1.0) circle (1.5pt) node[below right] {$P$};
			% forza risultante f
			\draw[->, thick] (1.5,1.0) -- (3.3,2.5)
			node[above right] {$\fb$};
			% componente f1 (orizzontale)
			\draw[->, thick] (1.5,1.0) -- (3.3,1.0)
			node[below] {$\fb_1$};
			% componente f2 (verticale)
			\draw[->, thick] (1.5,1.0) -- (1.5,2.5)
			node[left] {$\fb_2$};
			% parallelogramma tratteggiato
			\draw[dashed] (3.3,1.0) -- (3.3,2.5);
			\draw[dashed] (1.5,2.5) -- (3.3,2.5);
	\end{tikzpicture}}
	\caption{Operazioni elementari di equivalenza (I).
		\textbf{\textit{(a)}}~Traslazione della forza $\fb_k$
		lungo la propria linea d'azione da $P_k$ a $P_k'$:
		risultante e momento risultante restano invariati.
		\textbf{\textit{(b)}}~Composizione di due forze
		concorrenti in $P$: il sistema è equivalente alla
		risultante $\fb_1+\fb_2$ applicata in $P$.
		\textbf{\textit{(c)}}~Decomposizione della forza $\fb$
		applicata in $P$ nelle componenti $\fb_1$ e $\fb_2$
		secondo due direzioni indipendenti: operazione inversa
		alla composizione.}
	\label{fig:operazioni_elementari_1}
\end{figure}

\paragraph{Aggiunta o rimozione di un sistema equilibrato.}
In un punto qualsiasi $P$ si possono aggiungere (o rimuovere)
due forze uguali e opposte $\fb$ e $-\fb$ senza modificare gli
elementi caratteristici del sistema, poiché il loro contributo
alla risultante e al momento risultante è identicamente nullo.





\begin{oss}
	Le operazioni elementari possono essere combinate per
	ottenere trasformazioni più complesse. Ad esempio, il
	\textit{trasporto di una forza in un punto non appartenente\index{sistema di forze!trasporto di una forza}
		alla linea d'azione} si realizza come segue: data
	$(P,\,\fb)$, si aggiunge in un punto $Q$ la coppia
	equilibrata $(\fb,\,-\fb)$; si compone poi la forza
	originale $(P,\,\fb)$ con $(Q,\,-\fb)$, ottenendo una
	coppia di momento
	$\bm{\upmu} = \overrightarrow{QP}\times\fb$. Il sistema
	risultante è la forza $(Q,\,\fb)$ più la coppia
	$\bm{\upmu}$: la forza è stata ``trasportata'' in $Q$ al
	prezzo di una coppia correttiva (\FIGref{fig:trasporto_forza})
	\begin{figure}[htbp]
		\centering
		\subfloat[][\emph{sistema di partenza}]{%
			\begin{tikzpicture}[>=Stealth, line width=0.8pt, scale=0.80]
				\fill (0,0) circle (1.5pt) node[below] {$Q$};
				\fill (1.0,2.0) circle (1.5pt) node[above right] {$P$};
				% vettore QP
				\draw[->, thin, gray!60] (0,0) -- (1.0,2.0)
				node[midway, left] {$\overrightarrow{QP}$};
				% forza in P
				\draw[->, thick] (1.0,2.0) -- (2.8,2.0)
				node[above] {$\fb$};
		\end{tikzpicture}}
		\qquad
		\subfloat[][\emph{aggiunta della coppia equilibrata in $Q$}]{%
			\begin{tikzpicture}[>=Stealth, line width=0.8pt, scale=0.80]
				\fill (0,0) circle (1.5pt) node[below] {$Q$};
				\fill (1.0,2.0) circle (1.5pt) node[above right] {$P$};
				% vettore QP
				\draw[->, thin, gray!60] (0,0) -- (1.0,2.0)
				node[midway, left] {$\overrightarrow{QP}$};
				% forza originale in P
				\draw[->, thick] (1.0,2.0) -- (2.8,2.0)
				node[above] {$\fb$};
				% forza +f in Q
				\draw[->, thick] (0,0) -- (1.8,0)
				node[below] {$\fb$};
				% forza -f in Q
				\draw[->, thick] (0,0) -- (-1.8,0)
				node[below] {$-\fb$};
				% braccio b
				\draw[<->, thin] (2.9,0) -- (2.9,2.0)
				node[midway, right] {$b$};
		\end{tikzpicture}}
		\qquad
		\subfloat[][\emph{sistema equivalente}]{%
			\begin{tikzpicture}[>=Stealth, line width=0.8pt, scale=0.80]
				\fill (0,0) circle (1.5pt) node[below] {$Q$};
				\fill (1.0,2.0) circle (1.5pt) node[above right] {$P$};
				% forza risultante in Q
				\draw[->, thick] (0,0) -- (1.8,0)
				node[below] {$\fb$};
				% coppia correttiva (freccia semicircolare oraria)
				\draw[->, thick] (-0.6,0) arc[start angle=180,
				end angle=90, radius=0.6]
				node[above left] {$\mu(Q)$};
				% braccio b
				\draw[<->, thin] (2.9,0) -- (2.9,2.0)
				node[midway, right] {$b$};
		\end{tikzpicture}}
		\caption{Trasporto di una forza in un punto non appartenente
			alla linea d'azione.
			\textbf{\textit{(a)}}~Sistema originale: forza $\fb$
			applicata in $P$.
			\textbf{\textit{(b)}}~Aggiunta in $Q$ della coppia
			equilibrata $(\fb,\,-\fb)$: la forza $(P,\,\fb)$ e
			$(Q,\,-\fb)$ formano una coppia di momento
			$\bm{\upmu} = \protect\overrightarrow{QP}\times\fb$,
			con braccio $b$.
			\textbf{\textit{(c)}}~Sistema equivalente: forza $\fb$
			applicata in $Q$ più la coppia correttiva oraria di intensità 
			$\mu(Q) = |\fb|\,b$.}
		\label{fig:trasporto_forza}
	\end{figure}.
\end{oss}

\begin{oss}[\textbf{Riduzione di un sistema al polo}]
	Le operazioni elementari consentono di \textit{ridurre} un
	qualsiasi sistema di forze $\mathcal{S}$ a un polo $O$
	fissato. Si trasporta ciascuna forza $\fb_k$ nel polo $O$
	mediante aggiunta della coppia correttiva
	$\overrightarrow{OP_k}\times\fb_k$; si compongono poi tutte
	le forze concorrenti in $O$ nella risultante $\fb$ e tutte
	le coppie nel momento risultante $\bm{\upmu}(O)$. Il sistema
	così ottenuto --- la forza $\fb$ applicata in $O$ e la
	coppia $\bm{\upmu}(O)$ --- è equivalente al sistema di
	partenza e prende il nome di \textit{sistema ridotto al\index{sistema di forze!sistema ridotto al polo}
		polo~$O$}. Si noti che la riduzione dipende dalla scelta
	del polo: al variare di $O$ cambia il momento risultante
	secondo la~\eqref{eq:variazione_momento}, mentre la
	risultante resta invariata.
\end{oss}
















%
%
%
%
%
%\newpage
%\todo[inline]{Vecchia versione}
%%---------------------------------------------------------------------------------------------------------------------------- 
%\section{Legge di variazione del momento}\index{momento!legge di variazione|textbf}
%%---------------------------------------------------------------------------------------------------------------------------- 
%
%Il momento risultante dipende dal polo. Per determinare come
%varia, valutiamo la definizione~\eqref{eq:momento_risultante}
%in un polo $P$ qualsiasi. Poiché
%$\overrightarrow{PP_k} = \overrightarrow{OP_k} -
%\overrightarrow{OP}$:
%\begin{equation}
%	\bm{\upmu}(P)
%	= \sum_k \overrightarrow{PP_k} \times \fb_k
%	= \sum_k \overrightarrow{OP_k} \times \fb_k
%	- \overrightarrow{OP} \times \sum_k \fb_k
%	= \bm{\upmu}(O) - \overrightarrow{OP} \times \fb.
%\end{equation}
%Applicando l'antisimmetria del prodotto vettoriale:
%\begin{equation}\label{eq:variazione_momento}
%	\boxed{\bm{\upmu}(P)
%		= \bm{\upmu}(O) + \fb \times \overrightarrow{OP}.}
%\end{equation}
%Il momento non è un vettore libero: dipende dal polo
%attraverso un termine che coinvolge la risultante (\FIGref{fig:variazione_momento}). Conoscendo
%$\fb$ e $\bm{\upmu}(O)$, si può calcolare il momento rispetto
%a qualsiasi polo senza tornare alle singole forze.
%
%\begin{figure}[htbp]
%	\centering
%	% --- helper locali usati dai due pannelli ---
%	\providecommand{\ricciolo}[2]{%
%		\draw[->, thin, rotate around={#2:(#1)}]
%		([shift={(0.18,0)}]#1) arc[x radius=0.18, y radius=0.07,
%		start angle=0, end angle=305];}
%	\tikzset{midarrow/.style={postaction={decorate},
%			decoration={markings, mark=at position #1 with {\arrow{Stealth}}}},
%		midarrow/.default=0.55}
%	%
%	\subfloat[]{%
%		\begin{tikzpicture}[>=Stealth, line width=0.8pt, scale=1.0,
%			baseline={(current bounding box.center)}]
%			\coordinate (O)  at (0,0);
%			\coordinate (Pk) at (1.85,1.95);
%			\coordinate (P)  at (2.05,-0.35);
%			% triangolo dei vettori posizione (costruzioni)
%			\draw[thin, gray!85, midarrow=0.55] (O) -- (Pk);
%			\draw[thin, gray!85, midarrow=0.55] (O) -- (P);
%			\draw[thin, gray!85, midarrow=0.62] (P) -- (Pk);
%			%\node[gray!75] at ($(O)!0.5!(Pk)+(-0.32,0.42)$) {$\overrightarrow{OP_k}$};
%			
%			\draw[thin, gray!85] (O) -- (Pk) node[font=\small, midway, sloped, above] {$\overrightarrow{OP_k}$};
%			\draw[thin, gray!85] (O) -- (P) node[font=\small, midway, sloped, below] {$\overrightarrow{OP}$};
%			%	\draw[thin] (P) -- (Pk) node[midway, sloped, below] {$\overrightarrow{PP_k}$};
%			
%			%				\node[gray!75] at ($(O)!0.5!(P)+(0.05,-0.36)$)  {$\overrightarrow{OP}$};
%			\node[gray!75, font=\small] at ($(P)!0.5!(Pk)+(0.45,0.02)$) {$\overrightarrow{PP_k}$};
%			% forza f_k
%			\draw[->, thick] (Pk) -- ($(Pk)+(-0.50,1.05)$) node[right=1pt] {$\fb_k$};
%			% momento mu(O)
%			\coordinate (muOt) at (0,1.30);
%			\coordinate (muOc) at (0,1.52);
%			\draw[->>, thick] (O) -- (muOt) node[left=1pt] {$\bm{\upmu}(O)$};
%			%\ricciolo{muOc}{0}
%			% punti
%			\fill (O)  circle (1.4pt) node[font=\small, below left=-2pt]  {$O$};
%			\fill (P)  circle (1.4pt) node[font=\small, below right=-1pt] {$P$};
%			\fill (Pk) circle (1.4pt) node[font=\small, above right=-1pt] {$P_k$};
%	\end{tikzpicture}}
%	\qquad
%	\subfloat[]{%
%		\begin{tikzpicture}[>=Stealth, line width=0.8pt, scale=1.0,
%			baseline={(current bounding box.center)}]
%			\coordinate (O) at (0,0.55);
%			\coordinate (P) at (2.30,-0.45);
%			% OP (costruzione)
%			\draw[thin, gray!55, midarrow=0.55] (O) -- (P);
%			\node[gray!75, font=\small] at ($(O)!0.5!(P)+(0.08,0.34)$) {$\overrightarrow{OP}$};
%			% momento mu(O)
%			\coordinate (muOt) at (0,1.75);
%			\coordinate (muOc) at (0,1.97);
%			\draw[->>, thick] (O) -- (muOt) node[font=\small, left=1pt] {$\bm{\upmu}(O)$};
%			%\ricciolo{muOc}{0}
%			% risultante f in O
%			\draw[->, thick] ($(O)$) -- ($(O)+(0.92,1.00)$) node[right=-1pt] {$\fb$};
%			%\draw[->, thick] ($(O)+(0.70,0.02)$) -- ($(O)+(0.92,1.00)$) node[right=-1pt] {$\fb$};
%			% risultante f trasferita in P
%			\draw[->, thick] (P) -- ($(P)+(0.92,1.00)$) node[above right=-3pt] {$\fb$};
%			% momento mu(P) (assiale, obliquo)
%			\coordinate (muPt) at ($(P)+(0.98,-0.20)$);
%			\coordinate (muPc) at ($(P)+(1.16,-0.71)$);
%			\draw[->>, thick] (P) -- (muPt) node[right=1pt, yshift=3pt] {$\bm{\upmu}(P)$};
%			%\ricciolo{muPc}{-121}
%			% punti
%			\fill (O) circle (1.4pt) node[font=\small, left=2pt]        {$O$};
%			\fill (P) circle (1.4pt) node[font=\small, below left=-2pt] {$P$};
%			% legge di variazione
%			\node[anchor=west, font=\small] at (-0.2,-1.60)	{$\bm{\upmu}(P) = \bm{\upmu}(O) + \fb \times \overrightarrow{OP}$};
%	\end{tikzpicture}}
%	\caption{Rappresentazione grafica della legge di variazione
%		dei momenti: \textbf{\textit{(a)}}~calcolo del momento
%		prodotto da una forza $\fb_k$ rispetto a un polo di
%		riferimento $O$; \textbf{\textit{(b)}}~variazione del
%		momento risultante al cambiare del polo da $O$ a $P$,
%		dove il termine additivo $\fb \times
%		\protect\overrightarrow{OP}$ tiene conto dell'effetto
%		del trasferimento della risultante $\fb$ dal polo $O$
%		al polo $P$.}
%	\label{fig:variazione_momento}
%\end{figure}
%
%
%
%\begin{oss}
%	Se $\fb = \zerovb$ (coppia pura),
%	la~\eqref{eq:variazione_momento} dà
%	$\bm{\upmu}(P) = \bm{\upmu}(O)$ per ogni $P$: il momento è
%	indipendente dal polo, coerentemente con quanto già
%	osservato per la coppia.
%\end{oss}
%
%\begin{oss}
%	La struttura della~\eqref{eq:variazione_momento} è
%	formalmente identica alla formula fondamentale della
%	cinematica infinitesima del \CAPref{cap:cinematica_infinitesima}. Tuttavia, la
%	natura dei due risultati è diversa: nella cinematica, la
%	formula esprime direttamente il campo degli spostamenti ---
%	è la legge del campo, e i punti $O$ e $P$ hanno lo stesso
%	ruolo; nella statica, la~\eqref{eq:variazione_momento} è un
%	\textit{teorema} che si dimostra a partire dalla definizione
%	di momento risultante, e il polo ha un ruolo privilegiato.
%	Il dualismo cinematico-statico, che sarà formalizzato nella
%	Sezione~\ref{sec:dualismo}, risiede nel fatto che i due
%	campi obbediscono alla stessa legge affine, indipendentemente
%	da come ci si arriva.
%\end{oss}
%
%%---------------------------------------------------------------------------------------------------------------------------- 
%\section{Invarianti del campo dei momenti}	\label{sec:invarianti_F}
%%---------------------------------------------------------------------------------------------------------------------------- 
%
%Il campo dei momenti possiede due \textit{invarianti},
%grandezze indipendenti dalla scelta del polo~$O$:
%\begin{enumerate}[label=(\roman*)]
%	\item la risultante $\fb$, che compare
%	nella~\eqref{eq:variazione_momento} come dato del sistema
%	e non del polo;
%	\item il prodotto scalare $\fb \cdot \bm{\upmu}(O)$, che 
%	rappresenta la componente del momento nella direzione 
%	della risultante. Tale componente è la stessa per tutti 
%	i poli: questa proprietà prende il nome di 
%	\textit{equiproiettività}\index{equiproiettivita@equiproiettività} del campo dei momenti, e 
%	costituisce l'analogo statico dell'equiproiettività del 
%	campo degli spostamenti rigidi infinitesimi 
%	(cfr.~\CAPref{cap:cinematica_infinitesima}). Per 
%	verificarla è sufficiente proiettare 
%	la~\eqref{eq:variazione_momento} nella direzione di $\fb$:
%	\begin{equation}
%		\fb \cdot \bm{\upmu}(P)
%		= \fb \cdot \bm{\upmu}(O)
%		+ \fb \cdot (\fb \times \overrightarrow{OP})
%		= \fb \cdot \bm{\upmu}(O),
%	\end{equation}
%	essendo nullo il prodotto misto con due fattori uguali.
%\end{enumerate}
%
%\begin{oss}
%	Se $P$ si trova sulla retta per $O$ parallela a $\fb$
%	($\overrightarrow{OP} = \lambda\fb$), nella
%	legge di variazione il termine $\fb \times \lambda\fb$
%	si annulla identicamente. In tal caso non solo il prodotto
%	scalare, ma l'intero vettore momento è invariante:
%	$\bm{\upmu}(P) = \bm{\upmu}(O)$. Tutti i poli allineati
%	con la direzione della risultante danno lo stesso momento
%	risultante. Questa proprietà, più forte della semplice
%	invarianza scalare, è l'analogo statico dell'invarianza
%	dello spostamento per punti allineati con la direzione
%	della rotazione, già dimostrata nel \CAPref{cap:cinematica_infinitesima}.
%\end{oss}
%
%%---------------------------------------------------------------------------------------------------------------------------- 
%\section{Asse centrale del sistema di forze}
%%---------------------------------------------------------------------------------------------------------------------------- 
%
%Quando $\fb \neq \zerovb$, esiste una retta, l'\textit{asse
%	centrale del sistema di forze}\index{asse centrale!del sistema di forze|textbf}, lungo la quale il momento è
%parallelo alla risultante. La sua costruzione segue un
%ragionamento puramente geometrico, del tutto analogo a quello
%sviluppato nel \CAPref{cap:cinematica_infinitesima} per l'asse centrale degli
%spostamenti.
%
%\subsection{Decomposizione del momento}
%
%Fissiamo un polo $O$ e decomponiamo il momento risultante
%nella componente parallela e in quella ortogonale alla
%risultante:
%\begin{equation}\label{eq:decomp_momento}
%	\bm{\upmu}(O) = \bm{\upmu}_{\parallel}(O)
%	+ \bm{\upmu}_{\perp}(O),
%\end{equation}
%con
%\begin{equation}
%	\bm{\upmu}_{\parallel}(O)
%	= \frac{\fb \cdot \bm{\upmu}(O)}{|\fb|^2}\,\fb,
%	\qquad
%	\bm{\upmu}_{\perp}(O)
%	= \bm{\upmu}(O) - \bm{\upmu}_{\parallel}(O).
%\end{equation}
%
%\subsection{Invarianza della componente parallela}
%
%Nella legge di variazione~\eqref{eq:variazione_momento}, il
%termine $\fb \times \overrightarrow{OP}$ è \textit{sempre}
%ortogonale a $\fb$, qualunque sia il polo $P$. Valutando la
%legge in un polo $P$ qualsiasi, la componente del momento
%parallela a $\fb$ risulta
%\begin{equation}
%	\bm{\upmu}_{\parallel}(P)
%	= \bm{\upmu}_{\parallel}(O)
%	\qquad \forall\, O,\, P:
%\end{equation}
%la componente del momento nella direzione della risultante è
%la stessa per tutti i poli. Questo fornisce una lettura
%geometrica dell'invarianza del prodotto scalare
%$\fb \cdot \bm{\upmu}$ già dimostrata nella
%Sezione~\ref{sec:invarianti_F}.
%
%\subsection{Annullamento della componente ortogonale}
%
%Cerchiamo un punto $C$ in cui il momento sia puramente
%parallelo a $\fb$, cioè tale che
%$\bm{\upmu}_{\perp}(C) = \zerovb$. Valutando
%la~\eqref{eq:variazione_momento} in $C$ e separando la
%componente perpendicolare a $\fb$:
%\begin{equation}\label{eq:cond_annull_stat}
%	\bm{\upmu}_{\perp}(C)
%	= \bm{\upmu}_{\perp}(O)
%	+ \fb \times \overrightarrow{OC}
%	= \zerovb.
%\end{equation}
%Poiché $\fb \times \overrightarrow{OC}$ dipende solo dalla
%componente di $\overrightarrow{OC}$ perpendicolare a $\fb$
%--- la componente parallela non contribuisce al prodotto
%vettoriale --- poniamo
%$\overrightarrow{OC} = \overrightarrow{OC}_{\!\perp}
%+ \lambda\,\fb/|\fb|$
%e la condizione diventa
%\begin{equation}\label{eq:cond_annull_stat2}
%	\fb \times \overrightarrow{OC}_{\!\perp}
%	= -\bm{\upmu}_{\perp}(O).
%\end{equation}
%Il primo membro è ortogonale sia a $\fb$ sia a
%$\overrightarrow{OC}_{\!\perp}$; deve esserlo anche il
%secondo. Poiché $\bm{\upmu}_{\perp}(O)$ è già ortogonale a
%$\fb$ per costruzione, resta la condizione che
%$\overrightarrow{OC}_{\!\perp}$ sia ortogonale a
%$\bm{\upmu}_{\perp}(O)$: nel piano $\pi$ perpendicolare a $\fb$,
%la direzione di $\overrightarrow{OC}_{\!\perp}$ è
%univocamente determinata. Quanto al modulo, l'applicazione
%$\vb \mapsto \fb \times \vb$ ristretta a tale piano è una
%rotazione di $\pi/2$ composta con un riscalamento di fattore
%$|\fb|$, da cui
%\begin{equation}
%	|\overrightarrow{OC}_{\!\perp}|
%	= \frac{|\bm{\upmu}_{\perp}(O)|}{|\fb|}.
%\end{equation}
%Poiché $\lambda$ resta libero, il luogo dei punti $C$ è una
%retta parallela a $\fb$: l'asse centrale del sistema di forze
%(\FIGref{f:asse_centrale_forze}).
%
%\begin{figure}[htbp]
%	\centering
%	\usetikzlibrary{arrows.meta,calc,decorations.markings,angles,quotes}
%	% --- proiezione obliqua: (x=profondita', y=oriz, z=vert) -> schermo ---
%	\def\ax{0.5}\def\ay{0.30}
%	\newcommand{\xyz}[3]{{#2 + (#1)*\ax},{#3 + (#1)*\ay}}
%	\tikzset{midarrow/.style={postaction={decorate},decoration={markings, mark=at position 0.56 with {\arrow{Stealth}}}}}
%	\begin{tikzpicture}[scale=1.4, >=Stealth, line width=0.8pt]
%		
%		%==== coordinate 3D (x=profondita', y=oriz, z=vert) ====
%		\coordinate (Pi1) at (\xyz{-1.3}{-1.0}{0});   % era {-1.3}{-0.7}
%		\coordinate (Pi2) at (\xyz{-1.3}{ 4.5}{0});   % era {-1.3}{ 3.5}
%		\coordinate (Pi3) at (\xyz{ 2.7}{ 4.5}{0});   % era { 2.7}{ 3.5}
%		\coordinate (Pi4) at (\xyz{ 2.7}{-1.0}{0});   % era { 2.7}{ 3.5}→ refuso: era { 2.7}{ 3.5}
%		
%		% O e suo asse
%		\coordinate (O)     at (\xyz{1.6}{0.5}{0});
%		\coordinate (Otop)  at (\xyz{1.6}{0.5}{2.7});
%		\coordinate (Obot)  at (\xyz{1.6}{0.5}{-0.8});
%		\coordinate (ThOa)  at (\xyz{1.6}{0.5}{1.10});
%		\coordinate (ThOb)  at (\xyz{1.6}{0.5}{1.90});
%		\coordinate (curlO) at (\xyz{1.6}{0.5}{2.12});
%		\coordinate (uparO) at (\xyz{1.6}{0.5}{0.90});
%		
%		% u_perp(O): ora verso il BASSO-SINISTRA
%		\coordinate (uperpO) at (\xyz{0.76}{0.26}{0});  % era {0.2}{0.1}
%		
%		% P_pi davanti-destra (OP_pi va in basso-destra): base 0.2,2.3 (era 1.886,2.5)
%		\coordinate (Ppi)   at (\xyz{0.2}{2.3}{0});
%		\coordinate (Ctop)  at (\xyz{0.2}{2.3}{2.7});
%		\coordinate (Cbot)  at (\xyz{0.2}{2.3}{-0.8});
%		\coordinate (ThCa)  at (\xyz{0.2}{2.3}{1.10});
%		\coordinate (ThCb)  at (\xyz{0.2}{2.3}{1.90});
%		\coordinate (curlC) at (\xyz{0.2}{2.3}{2.12});
%		\coordinate (uparC) at (\xyz{0.2}{2.3}{0.90});
%		
%		% prodotto vettoriale = -u_perp(O): ora verso l'ALTO-DESTRA
%		\coordinate (cross)  at (\xyz{1.04}{2.54}{0});  % era {1.6}{2.7}
%		\coordinate (Bpar) at (\xyz{-0.64}{2.06}{0});
%		
%		
%		%==== piano pi ====
%		\draw[gray!60] (Pi1) -- (Pi2) -- (Pi3) -- (Pi4) -- cycle;
%		\node[gray!75] at ($(Pi1)+(0.58,0.14)$) {$\pi$};
%		
%		%==== assi verticali ====
%		\draw[gray!55, thin] (Obot) -- (Otop);
%		\draw[gray!55, thin] (Cbot) -- (Ctop);
%		
%		%==== parallelogramma di lati u_perp(O) e OP_pi (evidenziato) ====
%		\filldraw[gray!10, draw=gray!55, thin]	(O) -- (uperpO) -- (Bpar) -- (Ppi) -- cycle;
%		
%		%==== triangolo O--C_perp--(f x OC_perp) ====
%		\fill[gray!10] (O) -- (Ppi) -- (cross) -- cycle;
%		
%		% u_perp(O) traslato in P_pi: opposto e collineare al prodotto vettoriale
%		\draw[->>,  thin] (Ppi) -- (Bpar);
%		%\node[font=\small, below=4pt, right] at (Bpar) {$\ub_{\perp}(O)$};
%		
%		%==== segmento O--(punta di vartheta x OP_pi) e ampiezza di rotazione vartheta ====
%		\draw[draw=gray!75, thin] (O) -- (cross);
%		%\pic[draw=black, ->, thin, angle radius=10mm, angle eccentricity=1.25,	"$\vartheta$"{font=\small, inner sep=1pt}] {angle = Ppi--O--cross};
%		
%		%==== vettori rotazione assiali + riccioli (sopra la doppia freccia) ====
%		\draw[->, thick] (ThOa) -- (ThOb) node[right=1pt] {$\fb$};
%		\draw[->, thick] (ThCa) -- (ThCb) node[right=1pt] {$\fb$};
%		
%		%==== costruzione nel piano ====
%		%\draw[->, thin] (O) -- (Ppi);
%		%\draw[midarrow, thin] (O) -- (Ppi);   % era  \draw[->, thin] (O) -- (Ppi);
%		%\node[font=\small, below=1pt] at ($(O)!0.30!(Ppi)$) {$\overrightarrow{OC_\perp}$};   % era 0.55, below=2pt
%		
%		\draw[midarrow,thin] (O) -- (Ppi) node[midway, sloped, below,xshift=-12pt] {$\overrightarrow{OC_\perp}$};
%		
%		
%		\draw[->>, thick] (O) -- (uperpO);
%		\node[above=1pt] at (uperpO) {$\bm{\upmu}_{\perp}(O)$};
%		\draw[->>, thick] (O) -- (uparO);
%		\node[right=1pt] at ($(O)!0.6!(uparO)$) {$\bm{\upmu}_{\parallel}$};
%		
%		%==== a P_pi ====
%		\draw[->>, thick] (Ppi) -- (uparC);
%		\node[right=1pt] at ($(Ppi)!0.6!(uparC)$) {$\bm{\upmu}_{\parallel}$};
%		\draw[->>, thick] (Ppi) -- (cross);
%		\node[below right=-6pt] at (cross) {$\fb\times\overrightarrow{OC_\perp}$};
%		
%		
%		
%		%==== punti ====
%		\fill (O)   circle (1.5pt) node[above right=-1pt] {$O$};
%		\fill (Ppi) circle (1.5pt) node[below right=-4pt] {$C_\perp$};
%		
%		%==== annotazione asse centrale ====
%		\node[align=left, font=\footnotesize] (lab) at ($(Ctop)+(1.8,-0.3)$)	{asse centrale della\\ sollecitazione};
%		\draw[->, thin] ($(lab.north west)+(0,-0.18)$) -- ($(Ctop)!0.18!(ThCb)$);
%		
%	\end{tikzpicture}%
%	\caption{Costruzione geometrica dell'asse centrale del
%		sistema di forze. Il momento rispetto a $O$ si decompone
%		nella componente $\bm{\upmu}^O_{\parallel}$, parallela a
%		$\fb$ e invariante per tutti i poli, e nella componente
%		$\bm{\upmu}^O_{\perp}$, giacente nel piano $\pi$
%		ortogonale a $\fb$. Il punto $C_\perp$, intersezione
%		dell'asse centrale con $\pi$, è determinato dalla
%		condizione $\fb \times \protect\overrightarrow{OC_\perp}
%		= -\bm{\upmu}^O_{\perp}$; l'asse centrale è la retta
%		per $C_\perp$ parallela a $\fb$.}
%	\label{f:asse_centrale_forze}
%\end{figure}
%
%
%
%
%\subsection{Proprietà dell'asse centrale}
%
%
%
%%---------------------------------------------------------------------------------------------------------------------------- 
%\section{Classificazione dei sistemi equivalenti}\index{sistema di forze!classificazione|textbf}
%%---------------------------------------------------------------------------------------------------------------------------- 
%
%In base ai valori degli invarianti si distinguono i seguenti
%casi:
%\begin{enumerate}[label=(\roman*)]
%	\item \textit{Dinamo (o sistema avvitante)}\index{sistema di forze!dinamo (sistema avvitante)}:
%	$\fb \neq \zerovb$,
%	$\bm{\upmu}(O) \neq \zerovb$,
%	$\fb \cdot \bm{\upmu}(O) \neq 0$. Il momento ha sia
%	componente nel piano ortogonale a $\fb$ sia componente
%	lungo $\fb$. Sull'asse centrale il momento è puramente
%	parallelo a $\fb$: il sistema, visto da un polo sull'asse
%	centrale, appare come una \textit{forza più una coppia
%		parallela}.
%	\item \textit{Forza unica}\index{sistema di forze!forza unica}:
%	$\fb \neq \zerovb$ e
%	$\fb \cdot \bm{\upmu}(O) = 0$. L'asse centrale è il luogo
%	dei poli a momento nullo; il sistema è equivalente a una
%	singola forza $\fb$ applicata su tale retta.
%	\item \textit{Coppia pura}\index{coppia!pura}:
%	$\fb = \zerovb$, $\bm{\upmu}(O) \neq \zerovb$. Il momento
%	è indipendente dal polo; l'asse centrale non esiste (o,
%	equivalentemente, è all'infinito).
%	\item \textit{Sistema nullo}: $\fb = \zerovb$,
%	$\bm{\upmu}(O) = \zerovb$. Nessun effetto statico.
%\end{enumerate}
%
%\begin{oss}
%	La corrispondenza con la classificazione dei moti rigidi
%	del \CAPref{cap:cinematica_infinitesima} è immediata: la dinamo corrisponde alla
%	rototraslazione, la forza unica alla rotazione pura, la
%	coppia pura alla traslazione pura, il sistema nullo al
%	riposo. Si noti l'inversione apparente: nella cinematica è
%	la rotazione pura ad annullare l'invariante scalare
%	($\bm{\upvartheta} \cdot \ub(O) = 0$, spostamento nullo
%	sull'asse centrale); nella statica è la forza unica
%	($\fb \cdot \bm{\upmu}(O) = 0$, momento nullo sull'asse
%	centrale). In entrambi i casi, l'asse centrale è il luogo
%	dei punti in cui il campo si annulla.
%\end{oss}
%
%
%%---------------------------------------------------------------------------------------------------------------------------- 
%\section{Particolarizzazione al caso piano}	\label{sec:caso_piano_F}
%%---------------------------------------------------------------------------------------------------------------------------- 
%
%Per un sistema di forze nel piano $(x,y)$, la risultante
%giace nel piano e il momento risultante è diretto lungo
%$\ab_z$: si riduce a uno scalare con segno, $\mu_z(O)$\nomenclature[G]{$\mu_z(O)$}{momento risultante nel piano (scalare con segno), rispetto a $O$}.
%
%\subsection{Legge di variazione nel piano}
%
%La~\eqref{eq:variazione_momento} proiettata lungo $\ab_z$
%fornisce:
%\begin{equation}\label{eq:variazione_piano_vett}
%	\mu_z(P) = \mu_z(O)
%	+ (\fb \times \overrightarrow{OP}) \cdot \ab_z.
%\end{equation}
%Il termine
%$(\fb \times \overrightarrow{OP}) \cdot \ab_z$ è lo scalare
%che misura il momento della risultante $\fb$ rispetto al polo
%$P$: geometricamente, è il prodotto dell'intensità di $\fb$
%per la distanza con segno della retta d'azione dal polo.
%Sviluppando in componenti:
%\begin{equation}\label{eq:variazione_piano}
%	\mu_z(P) = \mu_z(O)
%	+ X\bigl(y(P)-y(O)\bigr)
%	- Y\bigl(x(P)-x(O)\bigr).
%\end{equation}
%
%
%
%\begin{oss}
%	Nel caso piano l'invariante
%	$\fb \cdot \bm{\upmu}(O)$ è identicamente nullo, poiché
%	la risultante giace nel piano e il momento è ortogonale ad
%	esso. Dei casi elencati nella classificazione, nel piano si
%	presentano dunque soltanto la forza unica
%	($\fb \neq \zerovb$, con retta d'azione nel piano finito),
%	la coppia pura ($\fb = \zerovb$) e il sistema nullo. La
%	situazione è l'analogo statico di quanto osservato nel \CAPref{cap:cinematica_infinitesima}  per la cinematica piana, dove l'invariante
%	$\bm{\upvartheta} \cdot \ub(O)$ è identicamente nullo e si
%	presentano solo rotazione pura, traslazione pura e riposo.
%\end{oss}
%
%\subsection{Asse di sollecitazione}\label{sec:asse_soll}
%
%Quando $\fb \neq \zerovb$, l'asse centrale del sistema di
%forze interseca il piano del moto in una retta, l'\textit{asse
%	di sollecitazione}\index{asse di sollecitazione|textbf}, lungo la quale il momento è nullo. Poiché
%nel piano $\fb \cdot \bm{\upmu}(O) = 0$, il sistema è sempre
%equivalente alla sola risultante $\fb$ applicata sull'asse di
%sollecitazione: il momento si annulla su un'intera retta, non
%in un singolo punto.
%
%Per determinare la posizione dell'asse di sollecitazione,
%cerchiamo i punti $C$ tali che $\mu_z(C) = 0$.
%Dalla~\eqref{eq:variazione_piano}:
%\begin{equation}\label{eq:cond_CdR}
%	X\bigl(y(C)-y(O)\bigr)
%	- Y\bigl(x(C)-x(O)\bigr)
%	= -\mu_z(O).
%\end{equation}
%Questa è una singola equazione scalare in due incognite: il
%luogo dei punti a momento nullo è una retta nel piano. La
%sua direzione è quella di $\fb$, poiché una componente di
%$\overrightarrow{OC}$ parallela a $\fb$ non contribuisce al
%prodotto vettoriale $\fb \times \overrightarrow{OC}$ e quindi
%non modifica il momento. La distanza di tale retta dal polo
%$O$ è
%\begin{equation}\label{eq:CdR_distanza}
%	|\overrightarrow{OC}_{\!\perp}|
%	= \frac{|\mu_z(O)|}{|\fb|},
%\end{equation}
%dove $\overrightarrow{OC}_{\!\perp}$ è la componente di
%$\overrightarrow{OC}$ perpendicolare a $\fb$. In forma
%vettoriale, il piede della perpendicolare da $O$ all'asse di
%sollecitazione  (\FIGref{fig:asse_sollecitazione}\,a) è
%\begin{equation}\label{eq:vettore_OC_F}
%	\overrightarrow{OC}_{\!\perp}
%	= -\frac{\mu_z(O)}{|\fb|^2}\,
%	(\ab_z \times \fb).
%\end{equation}
%
%\begin{figure}[htbp]
%	\centering
%	\begin{tikzpicture}[>=Stealth, line width=0.8pt, scale=1.0]
%		
%		\coordinate (O) at (0,0);
%		\coordinate (f) at ({1.5*cos(30)},{1.5*sin(30)});
%		% C: mu_z=1.5, |f|=1.5 -> |OC|=1.0
%		% C = (mu_z/|f|^2)*(fy,-fx) = (0.5, -0.866)
%		\coordinate (C) at (0.500, -0.866);
%		
%		% Retta perpendicolare a f tratteggiata con freccia verso infinito
%		\draw[dashed, gray!60]
%		({-0.3*cos(-60)},{-0.3*sin(-60)}) --
%		({2.2*cos(-60)}, {2.2*sin(-60)});
%		\draw[->, dashed, gray!60]
%		({2.2*cos(-60)},{2.2*sin(-60)}) --
%		({2.7*cos(-60)},{2.7*sin(-60)})
%		node[right=1pt] {\small$\fb\to \mathbf{0}$};
%		
%		% --- Sistema in O: forza + coppia ---
%		% Vettore f in O
%		\draw[->, thick] (O) -- (f) node[above right=1pt] {$\fb$};
%		% Coppia mu_z in O (freccia semicircolare antioraria)
%		\draw[<-, thick] (-0.55, 0) arc[start angle=180, end angle=60, radius=0.55];
%		\node[left=2pt] at (-0.55, 0.3) {$\mu_z(O)$};
%		
%		% Angolo retto
%		\draw[thin]
%		({0.15*cos(30)},{0.15*sin(30)}) --
%		({0.15*cos(30)+0.15*cos(-60)},{0.15*sin(30)+0.15*sin(-60)}) --
%		({0.15*cos(-60)},{0.15*sin(-60)});
%		
%		% Quota |OC| = mu_z/|f| (traslata a sinistra oltre la tratteggiata)
%		\draw[<->, thin]
%		(-0.433, -0.250) --
%		(0.067, -1.116)
%		node[midway, left=3pt] {\small$\dfrac{|\mu_z(O)|}{|\fb|}$};
%		
%		% Punto O e punto C
%		\fill (O) circle (2pt) node[below left=-4pt] {$O$};
%		\fill[white] (C) circle (3pt);
%		\draw[thick] (C) circle (3pt);
%		\node[below right=2pt] at (C) {$C$};
%		
%		% --- Sistema equivalente in C: solo forza ---
%		% Vettore f in C (traslato)
%		\draw[->, thick] (C) -- ++({1.5*cos(30)},{1.5*sin(30)})
%		node[above right=1pt] {$\fb$};
%		% Nessuna coppia (mu_z(C)=0)
%		%				\node[below left=2pt] at ({0.5+1.5*cos(30)/2},{-0.866+1.5*sin(30)/2})
%		%				{\small$\mu_z(C){=}0$};
%		
%		% Freccia tratteggiata che collega i due sistemi
%		\draw[dashed, gray!40] (O) -- (C);
%		
%	\end{tikzpicture}
%	\caption{Centro di riduzione\index{centro di riduzione|seealso{asse di sollecitazione}} nel caso piano. Il sistema equivalente
%		a una forza $\fb$ applicata in $O$ con momento risultante $\mu_z(O)$
%		si riduce alla sola forza $\fb$ applicata nel punto $C$ della
%		perpendicolare a $\fb$ passante per $O$, a distanza
%		$|\protect\overrightarrow{OC}| = |\mu_z(O)|/|\fb|$.
%		Al tendere di $|\fb|\to 0$ a parità di $\mu_z(O)$ il centro
%		si allontana all'infinito: una coppia pura non ammette centro
%		di riduzione.}
%	\label{fig:asse_sollecitazione}
%\end{figure}
%
%
%\paragraph{Coordinate dell'asse di sollecitazione.}
%Sviluppando la~\eqref{eq:vettore_OC_F} in componenti, il
%piede della perpendicolare da $O$ all'asse di sollecitazione  (\FIGref{fig:asse_sollecitazione}\,b) è
%\begin{equation}\label{eq:CdR_coord}
%	x(C) =  \frac{\mu_z(O)\,Y}{X^2+Y^2},
%	\qquad
%	y(C) =  - \frac{\mu_z(O)\,X}{X^2+Y^2}.
%\end{equation}
%L'asse di sollecitazione è la retta passante per $C$ con la
%direzione di $\fb$. Il sistema di forze è equivalente alla
%sola risultante $\fb$ applicata in un qualsiasi punto di
%tale retta.
%
%Il campo dei momenti si esprime in forma particolarmente
%semplice se si sceglie $C$ come polo. Poiché
%$\mu_z(C) = 0$, la legge di
%variazione~\eqref{eq:variazione_piano} con $C$ al posto di
%$O$ dà direttamente
%\begin{equation}\label{eq:momento_da_CdR}
%	\mu_z(P) = X\bigl(y(P) - y(C)\bigr)
%	- Y\bigl(x(P) - x(C)\bigr),
%\end{equation}
%che è il momento della sola risultante $\fb$ applicata in $C$
%rispetto al polo $P$.
%
%\begin{oss}
%	La posizione dell'asse di sollecitazione non dipende dal
%	polo $O$ scelto per la costruzione: se si valuta
%	la~\eqref{eq:variazione_piano} a partire da un diverso
%	polo $O'$, cambiano $\mu_z(O')$ e le coordinate di $O'$,
%	ma la retta su cui il momento si annulla resta la stessa.
%	Ciò è coerente con il fatto che l'asse centrale è un
%	invariante del sistema di forze.
%\end{oss}
%
%Nel caso di coppia pura ($\fb = \zerovb$,
%$\mu_z(O) \neq 0$), la~\eqref{eq:CdR_distanza} dà
%$|\overrightarrow{OC}| \to \infty$: l'asse di sollecitazione
%si allontana indefinitamente nella direzione perpendicolare a
%$\fb$.
%
%\begin{oss}
%	La coppia pura può essere riguardata come il caso limite
%	di una forza unica la cui intensità tende a zero mentre
%	il braccio si allontana all'infinito, in modo che il
%	prodotto $|\fb|\,|\overrightarrow{OC}| = |\mu_z(O)|$
%	resti finito. In termini di geometria proiettiva, l'asse
%	di sollecitazione è la retta impropria del piano: la
%	direzione della risultante è determinata, ma la retta
%	d'azione è all'infinito. Questa interpretazione è
%	l'analogo statico della traslazione pura nella cinematica,
%	dove il centro di rotazione è il punto improprio della
%	retta ortogonale allo spostamento.
%\end{oss}
%
%\subsection{Sistemi di forze equiorientate nel piano}\index{sistema di forze!equiorientate}
%
%Un caso frequente nella pratica strutturale è quello di forze
%dirette lungo una stessa direzione, applicate lungo un
%intervallo. La riduzione di questi sistemi è particolarmente
%semplice e consente di introdurre il concetto di
%\textit{baricentro delle forze}\index{baricentro delle forze|textbf}, che è l'analogo statico del
%baricentro dei centri di rotazione introdotto nel \CAPref{cap:cinematica_infinitesima}.
%
%\paragraph{Forze concentrate equiorientate.}
%Consideriamo $n$ forze $\fb_k = F_k\,\ab_y$ applicate nei
%punti di coordinata $x_k$ lungo un intervallo di lunghezza
%$\ell$. La risultante e il momento risultante rispetto
%all'origine $O$ sono
%\begin{equation}
%	\fb = \Bigl(\sum_{k=1}^{n} F_k\Bigr)\,\ab_y, \qquad
%	\bm{\upmu}(O) = \Bigl(\sum_{k=1}^{n} x_k\,F_k\Bigr)
%	\,\ab_x.
%\end{equation}
%L'asse di sollecitazione passa per il punto di coordinata
%\begin{equation}\label{eq:baricentro_forze}
%	d = \frac{\displaystyle\sum_{k=1}^{n} x_k\,F_k}
%	{\displaystyle\sum_{k=1}^{n} F_k},
%\end{equation}
%che è il baricentro dei punti di applicazione pesati con le
%rispettive intensità. Si noti l'analogia con
%la~\eqref{eq:comp_rot_baricentro} del \CAPref{cap:cinematica_infinitesima}: il
%centro della rotazione composta è il baricentro dei centri
%di rotazione pesati con le ampiezze. Il dualismo
%cinematico-statico opera qui con la corrispondenza
%$\vartheta_k \leftrightarrow F_k$,
%$C_k \leftrightarrow P_k$.
%
%\begin{esempio}[Forze concentrate su un intervallo]
%	Si consideri un intervallo di lunghezza $\ell$ soggetto a
%	forze concentrate $-F_k\,\ab_y$ con $F_k > 0$, dirette verso
%	il basso.
%	\begin{enumerate}[label=(\alph*)]
%		\item Due forze di uguale intensità $F$ applicate ai
%		terzi medi ($z_1 = \ell/3$, $z_2 = 2\ell/3$):
%		risultante $|\fb| = 2F$, punto di applicazione
%		$d = \ell/2$ (punto medio, per simmetria).
%		\item Due forze ai terzi medi con $F_1 = F$,
%		$F_2 = 2F$: risultante $|\fb| = 3F$, punto di
%		applicazione $d = 5\ell/9$, spostato verso la forza di
%		intensità maggiore.
%		\item Tre forze di uguale intensità $F$ agli estremi e
%		al punto medio ($z_1 = 0$, $z_2 = \ell/2$,
%		$z_3 = \ell$): risultante $|\fb| = 3F$, punto di
%		applicazione $d = \ell/2$ (per simmetria).
%	\end{enumerate}
%	\begin{figure}[htbp]
%		\centering
%		\subfloat[][\emph{due forze uguali ai terzi medi}]{%
%			\begin{tikzpicture}[>=Stealth, line width=0.8pt, scale=0.80]
%				\draw[thick] (0,0) -- (3,0);
%				\fill (0,0) circle (1.5pt) node[above] {$0$};
%				\fill (3,0) circle (1.5pt) node[above] {$\ell$};
%				% forze verso il basso
%				\fill (1,0) circle (1.5pt) node[above] {$\frac{\ell}{3}$};
%				\draw[->, thick] (1,0) -- (1,-1.2) node[left] {$F$};
%				\fill (2,0) circle (1.5pt) node[above] {$\frac{2\ell}{3}$};
%				\draw[->, thick] (2,0) -- (2,-1.2) node[right] {$F$};
%				% risultante verso il basso
%				\fill (1.5,0) circle (1.5pt);
%				\draw[->, thick, gray!50] (1.5,0) -- (1.5,-2.4)
%				node[left] {$2F$};
%				\draw[<->, thin] (0,1) -- (1.5,1)
%				node[midway, above] {$d=\frac{\ell}{2}$};
%		\end{tikzpicture}}
%		\qquad
%		\subfloat[][\emph{due forze diverse ai terzi medi}]{%
%			\begin{tikzpicture}[>=Stealth, line width=0.8pt, scale=0.80]
%				\draw[thick] (0,0) -- (3,0);
%				\fill (0,0) circle (1.5pt) node[above] {$0$};
%				\fill (3,0) circle (1.5pt) node[above] {$\ell$};
%				% forze verso il basso
%				\fill (1,0) circle (1.5pt) node[above] {$\frac{\ell}{3}$};
%				\draw[->, thick] (1,0) -- (1,-0.8) node[left] {$F$};
%				\fill (2,0) circle (1.5pt) node[above] {$\frac{2\ell}{3}$};
%				\draw[->, thick] (2,0) -- (2,-1.6) node[right] {$2F$};
%				% risultante verso il basso
%				\fill (1.667,0) circle (1.5pt);
%				\draw[->, thick, gray!50] (1.667,0) -- (1.667,-2.4)
%				node[left] {$3F$};
%				\draw[<->, thin] (0,1) -- (1.667,1)
%				node[midway, above] {$d=\frac{5\ell}{9}$};
%		\end{tikzpicture}}
%		\qquad
%		\subfloat[][\emph{tre forze uguali agli estremi e al centro}]{%
%			\begin{tikzpicture}[>=Stealth, line width=0.8pt, scale=0.80]
%				\draw[thick] (0,0) -- (3,0);
%				\fill (0,0) circle (1.5pt) node[above] {$0$};
%				\fill (3,0) circle (1.5pt) node[above] {$\ell$};
%				% forze verso il basso
%				\draw[->, thick] (0,0) -- (0,-1.2) node[left] {$F$};
%				\fill (1.5,0) circle (1.5pt) node[above] {$\frac{\ell}{2}$};
%				\draw[->, thick] (1.5,0) -- (1.5,-1.2) node[left] {$F$};
%				\draw[->, thick] (3,0) -- (3,-1.2) node[left] {$F$};
%				% risultante verso il basso
%				\draw[->, thick, gray!50] (1.5,0) -- (1.5,-2.4)
%				node[left] {$3F$};
%				\draw[<->, thin] (0,1) -- (1.5,1)
%				node[midway, above] {$d=\frac{\ell}{2}$};
%		\end{tikzpicture}}
%		\caption{Sistemi di forze concentrate $-F_k\,\ab_y$
%			equiorientate su un intervallo di lunghezza $\ell$:
%			la risultante (in grigio) è applicata nel baricentro
%			dei punti di applicazione pesati con le rispettive
%			intensità.
%			\textbf{\textit{(a)}}~$|\fb|=2F$, $d=\ell/2$.
%			\textbf{\textit{(b)}}~$|\fb|=3F$, $d=5\ell/9$.
%			\textbf{\textit{(c)}}~$|\fb|=3F$, $d=\ell/2$.}
%		\label{fig:forze_concentrate_intervallo}
%	\end{figure}
%\end{esempio}
%
%
%\begin{oss}
%	La composizione di forze parallele --- il cui punto di
%	applicazione risultante è il baricentro dei punti di
%	applicazione pesati con le intensità --- è l'analogo
%	statico della composizione di rotazioni infinitesime nel
%	piano, dove il centro risultante è il baricentro dei
%	centri pesati con le ampiezze
%	(\CAPref{cap:cinematica_infinitesima}). Il dualismo
%	cinematico-statico opera qui con la corrispondenza
%	$\vartheta_k \leftrightarrow F_k$, $C_k \leftrightarrow P_k$.
%\end{oss}
%
%
%\paragraph{Densità di forze equiorientate.}
%Se lungo l'intervallo agisce un carico distribuito
%$\fb_L(x) = p(x)\,\ab_y$, le sommatorie si sostituiscono
%con integrali:
%\begin{equation}
%	|\fb| = \int_0^{\ell} p(x)\dd x, \qquad
%	d = \frac{\displaystyle\int_0^{\ell}
%		x\,p(x)\dd x}{\displaystyle\int_0^{\ell}
%		p(x)\dd x}.
%\end{equation}
%L'intensità della risultante è l'area sottesa dal diagramma
%di carico $p(z)$; il punto di applicazione è il baricentro
%di tale diagramma (\FIGref{fig:carichi_distribuiti}\,a).
%
%I tre casi di maggiore interesse applicativo sono illustrati
%nella \FIGref{fig:carichi_distribuiti}\,b--d. Per un carico
%uniforme $p(x) = p$ la risultante vale $|\fb| = p\ell$ ed è
%applicata a $d = \ell/2$; per un carico linearmente variabile
%$p(x) = p\,x/\ell$ si ha $|\fb| = p\ell/2$ e $d = 2\ell/3$;
%per un carico parabolico $p(x) = p\,x^2/\ell^2$ si ha
%$|\fb| = p\ell/3$ e $d = 3\ell/4$. In tutti e tre i casi il
%punto di applicazione si sposta verso la zona di carico
%maggiore, e la sua posizione coincide con il baricentro della
%figura geometrica corrispondente --- rettangolo, triangolo e
%parabola.
%
%\begin{figure}[htbp]
%	\centering
%	\subfloat[][\emph{carico generico $p(x)$}]{%
%		\begin{tikzpicture}[>=Stealth, line width=0.8pt, scale=0.90]
%			% asse x
%			\draw[->] (0,0) -- (3.5,0) node[right] {$x$};
%			\node[below] at (0,0) {$0$};
%			\node[below] at (3,0) {$\ell$};
%			\draw[thick] (0,-0.1) -- (0,0.1);
%			\draw[thick] (3,-0.1) -- (3,0.1);
%			% curva generica: stessa funzione usata per le frecce
%			\draw[thick, domain=0:3, samples=100]
%			plot (\x, {1.0 + 0.6*sin(\x*120)});
%			\node[right] at (3, {1.0 + 0.6*sin(360)}) {$p(x)$};
%			% frecce che partono esattamente dalla curva
%			\foreach \x in {0.0, 0.2, 0.5, 0.8, 1.1, 1.4,
%				1.7, 2.0, 2.3, 2.55, 2.8, 3.0}{
%				\pgfmathsetmacro{\h}{1.0 + 0.6*sin(\x*120)}
%				\draw[->, thin] (\x, \h) -- (\x, 0);
%			}
%			% punto baricentro (calcolato numericamente ~1.5)
%			\fill (1.5,0) circle (1.5pt);
%			% risultante
%			\draw[->, thick] (1.5,0) -- (1.5,-1.2)
%			node[right] {$|\fb|$};
%			% quota d dal bordo sinistro al baricentro
%			\draw[<->, thin] (0,-0.5) -- (1.5,-0.5)
%			node[midway, below] {$d$};
%	\end{tikzpicture}}
%	\qquad
%	\subfloat[][\emph{carico uniforme}]{%
%		\begin{tikzpicture}[>=Stealth, line width=0.8pt, scale=0.90]
%			\draw[->] (0,0) -- (3.5,0) node[right] {$x$};
%			\node[below] at (0,0) {$0$};
%			\node[below] at (3,0) {$\ell$};
%			\draw[thick] (0,-0.1) -- (0,0.1);
%			\draw[thick] (3,-0.1) -- (3,0.1);
%			\foreach \x in {0, 0.375, 0.75, 1.125, 1.5,
%				1.875, 2.25, 2.625, 3.0}{
%				\draw[->, thin] (\x, 1.5) -- (\x, 0);
%			}
%			\draw[thick] (0,1.5) -- (3,1.5);
%			\node[left] at (0,1.5) {$p$};
%			\fill (1.5,0) circle (1.5pt);
%			\draw[->, thick] (1.5,0) -- (1.5,-1.2)
%			node[right] {$p\ell$};
%			\draw[<->, thin] (0,-0.5) -- (1.5,-0.5)
%			node[midway, below] {$\frac{\ell}{2}$};
%	\end{tikzpicture}}
%	\\[1em]
%	\subfloat[][\emph{carico lineare}]{%
%		\begin{tikzpicture}[>=Stealth, line width=0.8pt, scale=0.90]
%			\draw[->] (0,0) -- (3.5,0) node[right] {$x$};
%			\node[below] at (0,0) {$0$};
%			\node[below] at (3,0) {$\ell$};
%			\draw[thick] (0,-0.1) -- (0,0.1);
%			\draw[thick] (3,-0.1) -- (3,0.1);
%			\foreach \x in {0.375, 0.75, 1.125, 1.5,
%				1.875, 2.25, 2.625, 3.0}{
%				\draw[->, thin] (\x, {1.5*\x/3}) -- (\x, 0);
%			}
%			\draw[thick] (0,0) -- (3,1.5);
%			\node[right] at (3,1.5) {$p$};
%			\fill (2.0,0) circle (1.5pt);
%			\draw[->, thick] (2.0,0) -- (2.0,-1.2)
%			node[right] {$\frac{p\ell}{2}$};
%			\draw[<->, thin] (0,-0.5) -- (2.0,-0.5)
%			node[midway, below] {$\frac{2\ell}{3}$};
%	\end{tikzpicture}}
%	\qquad
%	\subfloat[][\emph{carico parabolico}]{%
%		\begin{tikzpicture}[>=Stealth, line width=0.8pt, scale=0.90]
%			\draw[->] (0,0) -- (3.5,0) node[right] {$x$};
%			\node[below] at (0,0) {$0$};
%			\node[below] at (3,0) {$\ell$};
%			\draw[thick] (0,-0.1) -- (0,0.1);
%			\draw[thick] (3,-0.1) -- (3,0.1);
%			\foreach \x in {0.375, 0.75, 1.125, 1.5,
%				1.875, 2.25, 2.625, 3.0}{
%				\draw[->, thin] (\x, {1.5*\x*\x/9}) -- (\x, 0);
%			}
%			\draw[thick, domain=0:3, samples=50]
%			plot (\x, {1.5*\x*\x/9});
%			\node[right] at (3,1.5) {$p$};
%			\fill (2.25,0) circle (1.5pt);
%			\draw[->, thick] (2.25,0) -- (2.25,-1.2)
%			node[right] {$\frac{p\ell}{3}$};
%			\draw[<->, thin] (0,-0.5) -- (2.25,-0.5)
%			node[midway, below] {$\frac{3\ell}{4}$};
%	\end{tikzpicture}}
%	\caption{Riduzione di carichi distribuiti
%		$\fb_L(x) = -p(x)\,\ab_y$ con $p(x) \geq 0$
%		su un intervallo di lunghezza $\ell$: la risultante
%		(freccia spessa) è applicata nel baricentro del
%		diagramma di carico.
%		\textbf{\textit{(a)}}~Caso generico $p(x)$:
%		risultante $|\fb|$ applicata a distanza $d$.
%		\textbf{\textit{(b)}}~Carico uniforme:
%		$|\fb|=p\ell$, $d=\ell/2$.
%		\textbf{\textit{(c)}}~Carico lineare:
%		$|\fb|=p\ell/2$, $d=2\ell/3$.
%		\textbf{\textit{(d)}}~Carico parabolico:
%		$|\fb|=p\ell/3$, $d=3\ell/4$.}
%	\label{fig:carichi_distribuiti}
%\end{figure}
%
%
%
%
%\newpage
%%-------------------------------------------------------------------------------
%\section{Esercizi proposti}
%%-------------------------------------------------------------------------------
%
%\begin{esercizio}[Riduzione di un sistema piano]
%	\label{ex:C02_riduzione_piano}
%	
%	Nel piano $(x,y)$, tre forze agiscono su un corpo rigido:
%	$\fb_1 = 2\,\ab_x$ applicata in $P_1 = (0,\,1)$,
%	$\fb_2 = -\ab_y$ applicata in $P_2 = (3,\,0)$,
%	$\fb_3 = \ab_x + 2\,\ab_y$ applicata in $P_3 = (1,\,2)$.
%	(a) Calcolare la risultante $\fb$.
%	(b) Calcolare il momento risultante $\mu(O)$ rispetto all'origine.
%	(c) Determinare un punto $C$ dell'asse di sollecitazione, ad esempio quello a minima distanza da $O$.
%	(d) Verificare che $\mu(C) = 0$ usando la~\eqref{eq:variazione_piano}.
%	
%	\risultato (a) $\fb = (3\;\,1)^\top$. (b) Sommando i momenti delle tre forze rispetto a $O$: $\mu_1 = 0 \cdot 0 - 1 \cdot 2 = -2$, $\mu_2 = 3 \cdot (-1) - 0 = -3$, $\mu_3 = 1 \cdot 2 - 2 \cdot 1 = 0$, da cui $\mu(O) = -5$. (c) L'asse di sollecitazione ha equazione $\mu(O) + X\,y - Y\,x = 0$, ovvero $-5 + 3y - x = 0$. Il punto a minima distanza da $O$ ha coordinate $C = (-0.5,\,1.5)$, con distanza $|\mu(O)|/|\fb| = 5/\sqrt{10}$. (d) $\mu(C) = \mu(O) + X\,y_C - Y\,x_C = -5 + 3 \cdot 1.5 - 1 \cdot (-0.5) = -5 + 4.5 + 0.5 = 0$.
%	
%\end{esercizio}
%
%\begin{esercizio}[Coppia di forze]
%	\label{ex:C02_coppia}
%	
%	Due forze $\fb_1 = 5\,\ab_y$ e $\fb_2 = -5\,\ab_y$ sono applicate rispettivamente in $P_1 = (1,\,0)$ e $P_2 = (4,\,0)$.
%	(a) Verificare che la risultante è nulla.
%	(b) Calcolare il momento della coppia rispetto a $O = (0,\,0)$.
%	(c) Verificare che il momento non cambia scegliendo un polo diverso, ad esempio $O' = (2,\,3)$.
%	
%	\risultato (a) $\fb = 5\,\ab_y - 5\,\ab_y = \zerovb$. (b) $\mu(O) = 1 \cdot 5 - 0 + 4 \cdot (-5) - 0 = 5 - 20 = -15$. Equivalentemente, $\bm{\upmu} = \overrightarrow{P_2 P_1} \times \fb_1 = (-3\;\,0\;\,0)^\top \times (0\;\,5\;\,0)^\top = (0\;\,0\;\,-15)^\top$. (c) Per la legge di variazione~\eqref{eq:variazione_piano} con $X = Y = 0$: $\mu(O') = \mu(O) = -15$. Verifica diretta rispetto a $O' = (2,\,3)$: $\mu_1 = (1-2)(5) - (0-3)(0) = -5$, $\mu_2 = (4-2)(-5) - (0-3)(0) = -10$, $\mu(O') = -15$. Il momento di una coppia è invariante rispetto al polo, come deve essere.
%	
%\end{esercizio}
%
%\begin{esercizio}[Legge di variazione del momento]
%	\label{ex:C02_variazione_momento}
%	
%	Per il sistema dell'\ESEref{ex:C02_riduzione_piano}:
%	(a) calcolare $\mu(O)$ direttamente dalla definizione, con polo $O = (0,\,0)$;
%	(b) calcolare $\mu(P)$ con $P = (2,\,1)$ direttamente dalla definizione;
%	(c) verificare che i risultati sono coerenti con la~\eqref{eq:variazione_piano}.
%	
%	\risultato (a) Già calcolato nell'\ESEref{ex:C02_riduzione_piano}: $\mu(O) = -5$. (b) Sommando i momenti rispetto a $P = (2,\,1)$: $\mu_1 = (0-2)(0) - (1-1)(2) = 0$, $\mu_2 = (3-2)(-1) - (0-1)(0) = -1$, $\mu_3 = (1-2)(2) - (2-1)(1) = -3$, da cui $\mu(P) = -4$. (c) Con la~\eqref{eq:variazione_piano}: $\mu(P) = \mu(O) + X(y_P - y_O) - Y(x_P - x_O) = -5 + 3 \cdot 1 - 1 \cdot 2 = -4$. I due valori coincidono.
%	
%\end{esercizio}
%
%\begin{esercizio}[Classificazione di sistemi nello spazio]
%	\label{ex:C02_classificazione_3D}
%	
%	Per ciascuno dei seguenti sistemi, calcolare gli invarianti $\fb$ e $\fb \cdot \bm{\upmu}(O)$ e classificare il sistema (sistema nullo, coppia pura, forza unica, dinamo):
%	(a) $\fb_1 = \ab_x + \ab_z$ applicata in $P_1 = (0,\,0,\,0)$ e $\fb_2 = -\ab_x - \ab_z$ applicata in $P_2 = (1,\,0,\,0)$;
%	(b) $\fb_1 = \ab_y$ applicata in $P_1 = (1,\,0,\,0)$ e $\fb_2 = 2\,\ab_y$ applicata in $P_2 = (0,\,1,\,0)$;
%	(c) $\fb_1 = \ab_x$ applicata in $P_1 = (0,\,1,\,0)$ e $\fb_2 = \ab_z$ applicata in $P_2 = (1,\,0,\,0)$.
%	
%	\risultato (a) $\fb = \zerovb$, $\bm{\upmu}(O) = \overrightarrow{OP_2} \times \fb_2 = (1\;\,0\;\,0)^\top \times (-1\;\,0\;\,-1)^\top = (0\;\,1\;\,0)^\top \neq \zerovb$: \textit{coppia pura}. (b) $\fb = (0\;\,3\;\,0)^\top$, $\bm{\upmu}(O) = (1\;\,0\;\,0)^\top \times (0\;\,1\;\,0)^\top + (0\;\,1\;\,0)^\top \times (0\;\,2\;\,0)^\top = (0\;\,0\;\,1)^\top$. $\fb \cdot \bm{\upmu}(O) = 0$: \textit{forza unica} (la risultante è non nulla e l'invariante scalare è zero). (c) $\fb = (1\;\,0\;\,1)^\top$, $\bm{\upmu}(O) = (0\;\,1\;\,0)^\top \times (1\;\,0\;\,0)^\top + (1\;\,0\;\,0)^\top \times (0\;\,0\;\,1)^\top = (0\;\,0\;\,-1)^\top + (0\;\,-1\;\,0)^\top = (0\;\,-1\;\,-1)^\top$. $\fb \cdot \bm{\upmu}(O) = -1 \neq 0$: \textit{dinamo}.
%	
%\end{esercizio}
%
%\begin{esercizio}[Trasporto di una forza]
%	\label{ex:C02_trasporto}
%	
%	Nel piano $(x,y)$, una forza $\fb = 3\,\ab_x + \ab_y$ è applicata in $P = (2,\,1)$.
%	(a) Utilizzando le operazioni elementari di equivalenza, trasportare la forza nell'origine $O = (0,\,0)$, determinando la coppia correttiva.
%	(b) Verificare che il sistema originale e quello ottenuto hanno la stessa risultante e lo stesso momento rispetto a $O$.
%	
%	\risultato (a) Trasportare $\fb$ da $P$ a $O$ aggiunge una coppia correttiva pari al momento di $\fb$ rispetto a $O$ cambiato di segno e applicato in $O$, ovvero $\mu_\mathrm{c} = -(x_P\,Y - y_P\,X) = -(2 \cdot 1 - 1 \cdot 3) = 1$. Il sistema equivalente è dunque: forza $\fb = (3\;\,1)^\top$ applicata in $O$, e coppia $\mu = 1$. (b) Sistema originale: risultante $(3\;\,1)^\top$, momento rispetto a $O$ pari a $\mu(O) = 2 \cdot 1 - 1 \cdot 3 = -1$. Sistema trasportato: risultante $(3\;\,1)^\top$, momento rispetto a $O$ pari a $0 + 1 = 1$. \emph{I momenti differiscono di segno}: la coppia correttiva va dunque scelta con segno opposto, $\mu_\mathrm{c} = +(x_P\,Y - y_P\,X) = -1$. Con questa scelta i due sistemi hanno entrambi risultante $(3\;\,1)^\top$ e momento $-1$ rispetto a $O$.
%	
%\end{esercizio}
%
%	
